\documentclass{iacrcc}
\usepackage[operators,sets]{cryptocode}
\usepackage{mathtools}
\usepackage{booktabs}
\usepackage{placeins}
% !TEX root = main.tex
\renewcommand{\vec}[1]{\mathbf{#1}}
\newcommand{\mat}[1]{\mathbf{#1}}
\newcommand{\Z}{\mathbb{Z}}
\newcommand{\R}{\mathbb{R}}
\newcommand{\ring}{\mathcal{R}}
\newcommand{\inner}[2]{\langle #1,#2\rangle}
\newcommand{\Te}{T_{\mathsf e}}
\newcommand{\To}{T_{\mathsf o}}
\newcommand{\Mold}{M_{\mathsf{orig}}}
\newcommand{\Unif}{\mathcal{U}}
\newcommand{\poly}{\mathsf{poly}}
\AtBeginDocument{\let\cref\Cref}
% Keep the class bibliography style and use numerical citation labels.
\makeatletter
\AtBeginDocument{%
  \renewcommand{\theHpclinenumber}{\arabic{Hpclinenumber}}%
  \renewcommand{\theH@pclinenumber}{\arabic{H@pclinenumber}}%
}
\AtBeginDocument{\renewcommand{\@lbibitem}[2][]{\@bibitem{#2}}}
\makeatother


\title[running={Optimal Acceptance for Iterative Rejection Sampling}]{Optimal Acceptance for Iterative Rejection Sampling}

\begin{document}
\maketitle

\begin{abstract}
Iterative rejection sampling constructs lattice signature responses through short signed shifts while preserving a prescribed Gaussian distribution. We determine acceptance functions with the smallest possible rejection probability for two moves and for four moves in two orthogonal directions, with matching lower bounds at every positive Gaussian width. The two-move functions fit directly into the existing framework. The four-move construction uses a norm-preserving ring rotation to supply its second direction. We give complete basic and compressed signature protocols, correctness proofs, and conditional unforgeability theorems. Across three explicit parameter sets, replacing the original two-move rule by four moves reduces estimated signature sizes from $1892$, $2557$, and $2991$ bytes to $1730$, $2362$, and $2732$ bytes at two expected sampler attempts. Keys and challenge distributions are fixed within each comparison. The reductions are $8.56\%$, $7.63\%$, and $8.66\%$. These are entropy-based estimates. We report the full lattice parameters and distinguish challenge-search contributions and lattice-cost estimates from the assumptions required for signature security.
\end{abstract}

\begin{textabstract}
Iterative rejection sampling constructs lattice signature responses through short signed shifts while preserving a prescribed Gaussian distribution. We determine acceptance functions with the smallest possible rejection probability for two moves and for four moves in two orthogonal directions, with matching lower bounds at every positive Gaussian width. The two-move functions fit directly into the existing framework. The four-move construction uses a norm-preserving ring rotation to supply its second direction. We give complete basic and compressed signature protocols, correctness proofs, and conditional unforgeability theorems. Across three explicit parameter sets, replacing the original two-move rule by four moves reduces estimated signature sizes from 1892, 2557, and 2991 bytes to 1730, 2362, and 2732 bytes at two expected sampler attempts. Keys and challenge distributions are fixed within each comparison. The reductions are 8.56%, 7.63%, and 8.66%. These are entropy-based estimates. We report the full lattice parameters and distinguish challenge-search contributions and lattice-cost estimates from the assumptions required for signature security.
\end{textabstract}

% !TEX root = main.tex
\section{Introduction}
\label{sec:intro}

Compact lattice signatures require short responses whose distribution hides the signing key. In Gaussian constructions, the response width affects both the encoded signature size and the probability that signing must restart. A wider Gaussian makes secret-dependent shifts easier to hide, but increases the signature length. Improving the acceptance rule at a fixed width can therefore reduce the number of sampling attempts or permit a narrower response at the same repetition rate.

Iterative rejection sampling~\cite{GartnerFull} constructs a response through a sequence of short signed shifts. Instead of correcting for the entire challenge-dependent shift at once, the signer processes one nonzero challenge coefficient at a time. Each step adds or subtracts a short vector, or rejects. Its acceptance probabilities are chosen so that an input Gaussian becomes the same Gaussian conditional on acceptance. The per-step repetition factors multiply over the challenge weight, making the choice of these probabilities relevant even when an individual step rarely rejects.

Lithium~\cite{cryptoeprint:2026/1790} studies the practical implementation of this approach through joint choices of parameters, Gaussian sampling, exponential evaluation, rejection sampling, and entropy coding. Its reported Lithium-120 implementation produces $1187$-byte signatures and signs in $166\,000$ cycles, compared with $2420$ bytes and $191\,000$ cycles for its ML-DSA-44 baseline. These results motivate studying the acceptance rule itself. Our question concerns the smallest rejection probability compatible with the prescribed Gaussian and allowed moves, rather than the implementation cost of evaluating a particular rule.

The framework of~\cite{GartnerFull} states three conditions on the two move probabilities and supplies functions satisfying them. Those conditions leave room to change the functions without changing the protocol. We determine the optimum within this framework. We also study four moves in two orthogonal directions, where a different parity constraint permits a smaller repetition factor.

\subsection{Contributions}

We give two results.
\begin{enumerate}
\item We construct acceptance functions by balancing the even and odd Gaussian masses along each shift direction. The construction weights the even terms of an alternating sum by the ratio of the two total masses. It gives one expression for each move probability, valid at every input, and fits directly into the three conditions of~\cite{GartnerFull}. For four moves, the same directional sums share one additional denominator. A norm-preserving ring rotation supplies the second orthogonal direction for every secret vector.
\item We prove that these functions minimise rejection for two moves, and for four moves in two equal-length orthogonal directions, at every positive Gaussian width. Each allowed move exchanges two parity classes. Their unequal masses impose a lower bound on rejection, which the functions attain. For four moves, the two one-dimensional mass imbalances multiply. We also identify the offsets that determine the common repetition factor on the discrete lattice.
\end{enumerate}

The two-move replacement preserves the keys, challenges, and verification equations of~\cite{GartnerFull}. At its three rejecting parameter sets, matching the original sampler repetition factors reduces the entropy-based signature estimates from $775$, $1184$, and $1694$ bytes to $764$, $1167$, and $1674$ bytes. Keeping the original widths instead reduces the expected sampler attempts by $27.5\%$, $41.4\%$, and $39.7\%$. These comparisons retain the original reduced challenge weights.

Four moves require a changed verification relation that retains half as many challenge bits as the ring degree. On three explicit parameter sets for this relation, we compare the rule of~\cite{GartnerFull}, optimal two moves, and optimal four moves, keeping the keys and challenge distribution fixed within each set. At two expected sampler attempts, four moves reduce the signature estimates from $1892$, $2557$, and $2991$ bytes to $1730$, $2362$, and $2732$ bytes, saving $8.56\%$, $7.63\%$, and $8.66\%$. The comparison is between acceptance rules on these parameters, not between the four-move scheme and the original degree-$256$ instantiations.

We give complete basic and compressed protocols with conditional unforgeability theorems in the classical random-oracle model. The four-move theorem uses a folded self-target assumption, whose relation is stated explicitly. Signature sizes are entropy estimates, and the reported repetition counts concern the iterative sampler before norm and encoding rejection. The parameter and impact tables appear in~\cref{sec:results}. Bit-dropping, full security proofs, and supporting parameter calculations are given in the appendices.

\subsection{Technical Overview}

\subsubsection{Mass Preservation and the Parity Bound}

Fix a lattice $L$, a Gaussian width $r>0$, and a nonzero shift $\vec v\in L$. Write $\rho_r(\vec y)=\exp(-\norm{\vec y}^2/(2r^2))$ for the unnormalised Gaussian mass, as formalised in~\cref{def:gaussian}. A step returns $\vec y-\vec v$ with probability $f_{\vec v}(\vec y)$ and $\vec y+\vec v$ with probability $g_{\vec v}(\vec y)$. To preserve the Gaussian with overall acceptance probability $1/M$, its incoming masses must satisfy
\[
 \rho_r(\vec y+\vec v)f_{\vec v}(\vec y+\vec v)
 +\rho_r(\vec y-\vec v)g_{\vec v}(\vec y-\vec v)
 =\rho_r(\vec y)/M.
\]
Together with nonnegative probabilities summing to at most one, this is the complete framework of~\cref{def:admissible}. The probability of accepting a particular input may vary with that input.

Every step remains on the line $\vec y+\Z\vec v$ and exchanges its even and odd indexed points. Let $E$ and $O$ be their respective total Gaussian masses, as in~\cref{def:parity}. Producing accepted even mass $E/M$ requires mass from odd inputs, whose total is only $O$. Thus $E/M\le O$. Exchanging the classes gives $O/M\le E$, so every admissible rule requires $M\ge\max\{E/O,O/E\}$. This lower bound depends only on the permitted moves, not on the form of the acceptance functions.

\subsubsection{Attaining the Bound}

The alternating sums used in~\cite{GartnerFull} suggest constructing probabilities from differences of neighbouring Gaussian masses. We first balance their totals. At an input $\vec y$, multiply every even term by $b_{\vec v}(\vec y)=O/E$, then subtract the adjacent odd term. The resulting differences sum to zero over the whole line. Their signs change once, from negative to positive, because the ratio of adjacent Gaussian masses is monotone. Every tail sum is consequently nonnegative. A tail starting after the sign change contains only nonnegative terms, while an earlier tail equals the negative of a nonpositive prefix.

The two directional tails, divided by $\rho_r(\vec y)M$, define $f_{\vec v}(\vec y)$ and $g_{\vec v}(\vec y)$ in~\cref{eq:sum,eq:functions}. Their sum is $b_{\vec v}(\vec y)/M$. The incoming tails cancel except for the required mass $\rho_r(\vec y)/M$. Hence the same class ratio both determines the probability budget and attains its lower bound. The largest imbalance occurs on the line through the origin. With $\alpha=r/\norm{\vec v}$, the exact common factor is
\[
 M^*(\alpha)=
 \frac{\sum_{j\in\Z}\exp(-(2j)^2/(2\alpha^2))}
 {\sum_{j\in\Z}\exp(-(2j+1)^2/(2\alpha^2))}.
\]
The origin belongs to the lattice, so it supplies an exact lower bound without a continuous approximation. The construction and optimality proof are given in~\cref{sec:construction}.

\subsubsection{A Second Orthogonal Direction}

With two orthogonal shifts of equal length, the Gaussian masses factor across the two integer coordinates. Every move exchanges points of even and odd total coordinate parity. The difference between these two class masses is the product of the one-dimensional differences. If $\delta_0=(M^*-1)/(M^*+1)$ is the largest normalised one-dimensional imbalance, the optimal two-move rejection is $2\delta_0/(1+\delta_0)$, whereas the four-move rejection is $2\delta_0^2/(1+\delta_0^2)$.

Simply selecting one of the two optimal directional samplers would retain the two-move repetition factor. Instead, divide both directional numerators by the same factor $1+b_1b_2$, where $b_i$ is the parity ratio in direction $i$. Moving in one direction inverts its ratio and leaves the other unchanged. This identity makes the incoming contributions add to the target Gaussian mass, while their total outgoing probability attains the two-dimensional parity bound. The resulting factor is
\[
 M_4^*(\alpha)=\frac12\left(M^*(\alpha)+\frac1{M^*(\alpha)}\right).
\]
The explicit probabilities and matching lower bound are given in~\cref{thm:four-optimal}.

\subsubsection{From Acceptance to Signature Size}

In the signature scheme, a binary challenge of Hamming weight $k$ requests $k$ short shifts derived from the secret vector $\vec s$. The public challenge is not centred. The sampler chooses a signed lift, and its mass-preservation identity gives the accepted Gaussian response exactly in the honest-verifier protocol. Thus neither a heuristic $\sqrt{k}$ expansion nor the worst-case bound $k\norm{\vec s}$ determines the response width. The worst-case challenge weight $\kappa$ determines the repetition exponent. A final retention test makes acceptance equal to $M^{-\kappa}$ for every permitted challenge, including variable-weight challenges.

For four moves in $\ring\cong\Z[X]/\langle X^d+1\rangle$, multiplication by $\omega=X^{d/2}$ rotates every coefficient vector to an orthogonal vector of the same norm. All four units $\{1,-1,\omega,-\omega\}$ become one modulo $\omega-1$. Verification therefore recovers the challenge through the parity map that adds the two halves of a ring element. This quotient has $d/2$ bits. Folding changes the commitment parity, not the dimension or norm of the response. The resulting loss of challenge capacity is included in the parameter choices, and the changed self-target relation remains explicit in~\cref{thm:unit-security}.

At fixed expected sampler attempts, a smaller optimal factor permits a smaller Gaussian width. Under the entropy model of~\cite{GartnerFull}, reducing the width from $r_{\mathsf{old}}$ to $r_{\mathsf{new}}$ saves $N\log_2(r_{\mathsf{old}}/r_{\mathsf{new}})$ bits before byte rounding, where $N$ is the full response dimension. The protocols impose a hard response-norm cutoff, and bit-dropping adds a worst-case reconstruction error. These bounds and the precise challenge distributions are retained in the size comparisons.

% !TEX root = main.tex
\section{The Rejection Sampling Framework}
\label{sec:framework}

\subsection{Gaussian Inputs and Signed Moves}

\begin{definition}[Lattice and Gaussian Distribution]
\label{def:gaussian}
Let $L\subset\R^n$ be a lattice, let $r>0$, and let $\vec v\in L\setminus\{\vec 0\}$. Vectors use their Euclidean inner product and norm. For $\vec y\in\R^n$ and a countable set $A\subset\R^n$, define
\[
\rho_r(\vec y)\coloneqq\exp\!\left(-\frac{\norm{\vec y}^2}{2r^2}\right),
\qquad
\rho_r(A)\coloneqq\sum_{\vec y\in A}\rho_r(\vec y).
\]
The centred discrete Gaussian $D_{L,r}$ assigns probability $\rho_r(\vec y)/\rho_r(L)$ to each $\vec y\in L$. Define the width-to-shift ratio $\alpha\coloneqq r/\norm{\vec v}$.
\end{definition}

This includes the coefficient lattice used in~\cite{GartnerFull}. For a power of two $d$, a positive integer $\ell$, and $\ring\cong\Z[X]/\langle X^d+1\rangle$, coefficient vectors identify $\ring^\ell$ with $\Z^{d\ell}$. Thus the Gaussian calculations below apply directly to its ring vectors.

\begin{definition}[Rejection Step]
\label{def:step}
For functions $f_{\vec v},g_{\vec v}\colon L\to\R$ satisfying $f_{\vec v}(\vec y),g_{\vec v}(\vec y)\ge0$ and $f_{\vec v}(\vec y)+g_{\vec v}(\vec y)\le1$, let $R_{\vec v}(\vec y)$ return $\vec y-\vec v$ with probability $f_{\vec v}(\vec y)$, return $\vec y+\vec v$ with probability $g_{\vec v}(\vec y)$, and return the failure symbol $\bot$ otherwise. Its procedure is given in~\cref{fig:step}, where $\Unif([0,1))$ denotes the uniform distribution on $[0,1)$.
\end{definition}

\begin{figure}[t]
\centering
\procedure[linenumbering]{$R_{\vec v}(\vec y)$}{
 a \coloneqq f_{\vec v}(\vec y) \\
 b \coloneqq g_{\vec v}(\vec y) \\
 u \leftarrow \Unif([0,1)) \\
 \textbf{if }u<a \\
 \quad\textbf{return }\vec y-\vec v \\
 \textbf{if }u<a+b \\
 \quad\textbf{return }\vec y+\vec v \\
 \textbf{return }\bot
}
\caption{The rejection step of~\cite{GartnerFull}.}
\label{fig:step}
\end{figure}

\begin{definition}[Admissible Acceptance Functions]
\label{def:admissible}
A pair $f_{\vec v},g_{\vec v}$ is admissible with repetition factor $M\ge1$ if, for every $\vec y\in L$,
\begin{align}
 f_{\vec v}(\vec y)+g_{\vec v}(\vec y)&\le1,\label{eq:budget}\\
 f_{\vec v}(\vec y),\ g_{\vec v}(\vec y)&\ge0,\label{eq:positive}\\
 \rho_r(\vec y+\vec v)f_{\vec v}(\vec y+\vec v)+\rho_r(\vec y-\vec v)g_{\vec v}(\vec y-\vec v)&=\frac{\rho_r(\vec y)}M.\label{eq:balance}
\end{align}
The same $M$ is used for all inputs.
\end{definition}

\begin{lemma}[Output Distribution, {\cite[Lemma~1]{GartnerFull}}]
\label{lem:output}
For admissible acceptance functions and $\vec y\leftarrow D_{L,r}$, let $\vec z\leftarrow R_{\vec v}(\vec y)$. Then
\[
 \Pr[\vec z=\vec x]=\frac{D_{L,r}(\vec x)}M\quad(\vec x\in L),
 \qquad
 \Pr[\vec z\ne\bot]=\frac1M,
 \qquad
 (\vec z\mid\vec z\ne\bot)\sim D_{L,r}.
\]
\end{lemma}

The two terms in~\cref{eq:balance} are the only ways to produce a given output. The normalising constant $\rho_r(L)$ is common to the input and target distributions and cancels. These are exactly the three conditions of~\cite{GartnerFull}, specialised to equal Gaussian input and output distributions.

\subsection{Gaussian Masses Along the Shift Direction}

\begin{definition}[Even and Odd Masses]
\label{def:parity}
For $\vec y\in L$, define
\[
 E_{\vec v}(\vec y)\coloneqq\sum_{j\in\Z}\rho_r(\vec y+2j\vec v),
 \qquad
 O_{\vec v}(\vec y)\coloneqq\sum_{j\in\Z}\rho_r(\vec y+(2j+1)\vec v),
 \qquad
 b_{\vec v}(\vec y)\coloneqq\frac{O_{\vec v}(\vec y)}{E_{\vec v}(\vec y)}.
\]
Both masses are positive. The even class contains the input $\vec y$.
\end{definition}

Only ratios of these Gaussian masses enter the construction. Their dependence on the ambient vector reduces to one scalar.

\begin{lemma}[One-Dimensional Gaussian Ratios]
\label{lem:line}
Let $t=\inner{\vec y}{\vec v}/\norm{\vec v}^2$. For every integer $k$,
\[
 \frac{\rho_r(\vec y+k\vec v)}{\rho_r(\vec y)}
 =\exp\!\left(-\frac{k^2+2kt}{2\alpha^2}\right).
\]
Moreover,
\[
 b_{\vec v}(\vec y+\vec v)=b_{\vec v}(\vec y-\vec v)=\frac1{b_{\vec v}(\vec y)},
 \qquad b_{\vec v}(-\vec y)=b_{\vec v}(\vec y).
\]
\end{lemma}
\begin{proof}
Expanding $\norm{\vec y+k\vec v}^2$ gives $\norm{\vec y}^2+2k\inner{\vec y}{\vec v}+k^2\norm{\vec v}^2$. Substitution into the Gaussian ratio proves the first identity. A shift by $\vec v$ or $-\vec v$ exchanges the two sums of~\cref{def:parity}. Reflection preserves each sum because $\rho_r(-\vec x)=\rho_r(\vec x)$ and the indices can be negated.
\end{proof}

% !TEX root = main.tex
\section{Optimal Two-Move Acceptance}
\label{sec:construction}

The alternating sum of~\cite{GartnerFull} subtracts odd Gaussian terms from even ones. We weight the even terms by $b_{\vec v}(\vec y)$. The total weighted even mass then equals the total odd mass, which makes both directional tail sums nonnegative.

\subsection{The Weighted Acceptance Functions}

\begin{definition}[Weighted Alternating Sum]
\label{def:sum}
For $\vec y\in L$, define
\begin{equation}
\label{eq:sum}
 \widehat S_{\vec v}(\vec y)\coloneqq
 \frac1{\rho_r(\vec y)}\sum_{j=0}^{\infty}
 \bigl(b_{\vec v}(\vec y)\rho_r(\vec y+2j\vec v)-\rho_r(\vec y+(2j+1)\vec v)\bigr).
\end{equation}
For $M\ge1$, the acceptance functions are
\begin{equation}
\label{eq:functions}
 \boxed{\displaystyle
 f_{\vec v}(\vec y)\coloneqq\frac{\widehat S_{\vec v}(\vec y)}M,
 \qquad
 g_{\vec v}(\vec y)\coloneqq\frac{\widehat S_{\vec v}(-\vec y)}M.}
\end{equation}
All sums are absolutely convergent.
\end{definition}

\subsection{Nonnegativity and Mass Preservation}

\begin{lemma}[Nonnegative Tail Sums]
\label{lem:tails}
Let $(a_j)_{j\in\Z}$ be absolutely summable, with $\sum_j a_j=0$. Suppose an integer $J_0$ satisfies $a_j\le0$ for $j<J_0$ and $a_j\ge0$ for $j\ge J_0$. Then $\sum_{j\ge J}a_j\ge0$ for every integer $J$.
\end{lemma}
\begin{proof}
For $J\ge J_0$, every term in the tail is nonnegative. For $J<J_0$, the tail equals $-\sum_{j<J}a_j$, which is nonnegative because every term in this prefix is nonpositive.
\end{proof}

\begin{theorem}[Correctness of the Acceptance Functions]
\label{thm:construction}
For every $\vec y\in L$, the functions of~\cref{def:sum} are nonnegative and satisfy~\cref{eq:balance}. Their sum is
\begin{equation}
\label{eq:sum-budget}
 f_{\vec v}(\vec y)+g_{\vec v}(\vec y)=\frac{b_{\vec v}(\vec y)}M.
\end{equation}
Consequently, they are admissible whenever $M\ge\sup_{\vec y\in L}b_{\vec v}(\vec y)$.
\end{theorem}
\begin{proof}
Fix $\vec y$ and abbreviate $p_k=\rho_r(\vec y+k\vec v)$, $E=E_{\vec v}(\vec y)$, $O=O_{\vec v}(\vec y)$, and $b=O/E$. Define $a_j=bp_{2j}-p_{2j+1}$ for all $j\in\Z$. The sequence is absolutely summable, and
\[
 \sum_{j\in\Z}a_j=bE-O=0.
\]
By~\cref{lem:line}, with $t=\inner{\vec y}{\vec v}/\norm{\vec v}^2$,
\[
 \frac{p_{2j+1}}{p_{2j}}=\exp\!\left(-\frac{2t+4j+1}{2\alpha^2}\right).
\]
This ratio strictly decreases from infinity to zero as $j$ increases. Since $b>0$, the sign of $a_j$ changes once from negative to positive. By~\cref{lem:tails}, $\sum_{j\ge0}a_j\ge0$, proving $\widehat S_{\vec v}(\vec y)\ge0$. Applying the same argument to $-\vec y$ proves $\widehat S_{\vec v}(-\vec y)\ge0$.

Reflection expresses $p_0\widehat S_{\vec v}(-\vec y)$ as $\sum_{j\ge0}(bp_{-2j}-p_{-(2j+1)})$. The two even half sums count $p_0$ twice and every other even term once. The odd half sums count each odd term once. Therefore
\[
 p_0\bigl(\widehat S_{\vec v}(\vec y)+\widehat S_{\vec v}(-\vec y)\bigr)
 =b(E+p_0)-O=bp_0,
\]
which proves~\cref{eq:sum-budget}.

For the incoming mass, the ratio at each of $\vec y+\vec v$ and $\vec y-\vec v$ is $1/b$. Expanding the two functions gives
\begin{align*}
 Mp_1 f_{\vec v}(\vec y+\vec v)
 &=\sum_{j\ge0}\bigl(b^{-1}p_{2j+1}-p_{2j+2}\bigr),\\
 Mp_{-1}g_{\vec v}(\vec y-\vec v)
 &=\sum_{j\ge0}\bigl(b^{-1}p_{-(2j+1)}-p_{-(2j+2)}\bigr).
\end{align*}
Their sum is $b^{-1}O-(E-p_0)=p_0$, proving~\cref{eq:balance}. Finally, the assumed bound on $M$ and~\cref{eq:sum-budget} give~\cref{eq:budget}.
\end{proof}

\subsection{Use in an Iterative Sampler}

The construction uses $M$ as a common parameter and computes $b_{\vec v}(\vec y)$ from the Gaussian masses at the input. By~\cref{lem:line}, evaluating the functions requires only $r$, $\norm{\vec v}^2$, and $\inner{\vec y}{\vec v}$.

\begin{corollary}[Fixed Sequences of Shifts]
\label{cor:iteration}
Let $\vec v_1,\ldots,\vec v_k\in L\setminus\{\vec0\}$ be fixed before sampling $\vec y\leftarrow D_{L,r}$. Suppose the functions for each shift are admissible with the same factor $M$. Starting with $\vec z_0=\vec y$, apply $\vec z_i\leftarrow R_{\vec v_i}(\vec z_{i-1})$ for $i=1,\ldots,k$, aborting on failure. Then
\[
 \Pr[\vec z_k=\vec x]=M^{-k}D_{L,r}(\vec x)
 \quad(\vec x\in L).
\]
For an integer $\kappa\ge k$, independently retaining a successful result with probability $M^{k-\kappa}$ makes the total acceptance probability $M^{-\kappa}$ and leaves the conditional output distribution equal to $D_{L,r}$.
\end{corollary}
\begin{proof}
The unnormalised distribution before step $i$ is $M^{-(i-1)}D_{L,r}$ by induction, starting with $i=1$. Applying~\cref{lem:output} multiplies it by $1/M$. This proves the displayed identity at $i=k$. Independent retention multiplies each output probability by $M^{k-\kappa}$, giving $M^{-\kappa}D_{L,r}(\vec x)$.
\end{proof}

The retention factor in~\cref{cor:iteration} is the one used in the iterative sampler of~\cite{GartnerFull}. The new functions supply its single-step Gaussian identity with the same two signed moves. A common parameter valid for every shift of a prescribed maximum length is given in~\cref{cor:bound}.

% !TEX root = main.tex
\subsection{Optimality}
\label{sec:optimality}

\subsubsection{The Parity Lower Bound}

The accepted even mass must come entirely from odd inputs. This gives a lower bound on $M$ that applies to every admissible pair of functions, including functions unrelated to~\cref{eq:sum}.

\begin{lemma}[Necessary Repetition Factor]
\label{lem:lower}
Every pair of admissible acceptance functions satisfies
\[
 M\ge\max\!\left\{\frac{E_{\vec v}(\vec y)}{O_{\vec v}(\vec y)},\frac{O_{\vec v}(\vec y)}{E_{\vec v}(\vec y)}\right\}
 \qquad\text{for every }\vec y\in L.
\]
\end{lemma}
\begin{proof}
Sum~\cref{eq:balance} over the points $\vec y+2j\vec v$, with $j\in\Z$. Each odd input contributes its mass multiplied by the sum of its two move probabilities. These probabilities sum to at most one, so
\[
 \frac{E_{\vec v}(\vec y)}M
 =\sum_{j\in\Z}\rho_r(\vec y+(2j+1)\vec v)
 \bigl(f_{\vec v}(\vec y+(2j+1)\vec v)+g_{\vec v}(\vec y+(2j+1)\vec v)\bigr)
 \le O_{\vec v}(\vec y).
\]
Summing over the odd outputs instead gives $O_{\vec v}(\vec y)/M\le E_{\vec v}(\vec y)$. All terms are nonnegative, which justifies rearranging the sums.
\end{proof}

\begin{corollary}[Optimal Factor on a Single Line]
\label{cor:line}
On the set $\vec y+\Z\vec v$, the smallest factor satisfying~\cref{eq:budget,eq:positive,eq:balance} is
\[
 \max\!\left\{E_{\vec v}(\vec y)/O_{\vec v}(\vec y),O_{\vec v}(\vec y)/E_{\vec v}(\vec y)\right\}.
\]
The functions of~\cref{def:sum} attain it.
\end{corollary}
\begin{proof}
Along this set, $b_{\vec v}$ alternates between $O_{\vec v}(\vec y)/E_{\vec v}(\vec y)$ and its reciprocal. Thus~\cref{thm:construction} attains the bound of~\cref{lem:lower}.
\end{proof}

\subsubsection{The Largest Imbalance}

\begin{definition}[Scalar Parity Masses]
\label{def:scalar}
For $\alpha>0$ and $t\in\R$, define
\[
 \Te(t)\coloneqq\sum_{j\in\Z}e^{-(t+2j)^2/(2\alpha^2)},
 \qquad
 \To(t)\coloneqq\sum_{j\in\Z}e^{-(t+2j+1)^2/(2\alpha^2)}.
\]
The dependence on $\alpha$ is implicit. Define
\begin{equation}
\label{eq:optimum}
 M^*(\alpha)\coloneqq\frac{\Te(0)}{\To(0)}
 =\frac{\sum_{j\in\Z}e^{-(2j)^2/(2\alpha^2)}}{\sum_{j\in\Z}e^{-(2j+1)^2/(2\alpha^2)}}.
\end{equation}
\end{definition}

\begin{lemma}[Centred Line Has the Largest Imbalance]
\label{lem:worst}
For every $\alpha>0$, one has $M^*(\alpha)>1$ and
\[
 \max_{t\in\R}\max\!\left\{\frac{\Te(t)}{\To(t)},\frac{\To(t)}{\Te(t)}\right\}=M^*(\alpha).
\]
Furthermore, $M^*(\alpha)$ strictly decreases with $\alpha$.
\end{lemma}
The proof is given in~\cref{app:analytic}.

\begin{theorem}[Optimal Acceptance on the Discrete Lattice]
\label{thm:optimal}
Let $L\subset\R^n$ be a lattice, $r>0$, and $\vec v\in L\setminus\{\vec0\}$. With $\alpha=r/\norm{\vec v}$, the functions of~\cref{def:sum} are admissible with $M=M^*(\alpha)$. This is the smallest repetition factor among all functions satisfying~\cref{eq:budget,eq:positive,eq:balance}. The minimum rejection probability is
\[
 1-\frac1{M^*(\alpha)}=1-\frac{\To(0)}{\Te(0)}.
\]
\end{theorem}
\begin{proof}
Fix $\vec y\in L$ and set $t=\inner{\vec y}{\vec v}/\norm{\vec v}^2$. The orthogonal decomposition $\vec y=\vec y_{\perp}+t\vec v$ yields
\[
 \rho_r(\vec y+k\vec v)=\exp\!\left(-\frac{\norm{\vec y_{\perp}}^2}{2r^2}\right)
 \exp\!\left(-\frac{(t+k)^2}{2\alpha^2}\right).
\]
The common first factor cancels from the parity ratio, so $b_{\vec v}(\vec y)=\To(t)/\Te(t)\le M^*(\alpha)$ by~\cref{lem:worst}. Admissibility follows from~\cref{thm:construction}.

For necessity, apply~\cref{lem:lower} at $\vec y=\vec0$, which is a point of $L$. The masses on $\Z\vec v$ are exactly $e^{-k^2/(2\alpha^2)}$, so $E_{\vec v}(\vec0)/O_{\vec v}(\vec0)=M^*(\alpha)$. Every admissible pair therefore requires $M\ge M^*(\alpha)$. The rejection probability follows from~\cref{lem:output}.
\end{proof}

\begin{corollary}[A Common Bound for Several Shifts]
\label{cor:bound}
Let $B>0$ bound every permitted shift length. The choice $M=M^*(r/B)$ makes the functions of~\cref{def:sum} admissible for all $\vec v\in L$ with $0<\norm{\vec v}\le B$. If a permitted shift has length $B$, this common factor is optimal.
\end{corollary}
\begin{proof}
Since $r/\norm{\vec v}\ge r/B$, monotonicity in~\cref{lem:worst} gives $M^*(r/\norm{\vec v})\le M^*(r/B)$. Apply~\cref{thm:construction,thm:optimal}. A permitted shift of length $B$ also gives the matching lower bound by~\cref{thm:optimal}.
\end{proof}

\subsection{Comparison with the Original Factor}

\begin{definition}[Original Repetition Factor, {\cite[Corollary~1]{GartnerFull}}]
\label{def:original}
The Gaussian repetition factor used in~\cite{GartnerFull} is
\begin{equation}
\label{eq:original}
 \Mold(\alpha)\coloneqq1+\frac{2\alpha\sqrt{2\pi}\,e^{-\pi^2\alpha^2/2}}{e^{-1/(2\alpha^2)}(1-e^{-2\pi^2\alpha^2})}.
\end{equation}
\end{definition}

\begin{proposition}[Strict Improvement]
\label{prop:improvement}
For every $\alpha>0$, $M^*(\alpha)<\Mold(\alpha)$. The acceptance functions of~\cite[Theorem~1]{GartnerFull} also require a repetition factor strictly larger than $M^*(\alpha)$ when their factor is optimised.
\end{proposition}
The proof is given in~\cref{app:analytic}.

% !TEX root = main.tex
\section{The Signature Scheme}
\label{sec:protocol}

We instantiate the acceptance functions in the signature protocol of~\cite{GartnerFull}. The basic protocol is given here, and its bit-dropping variant is specified in~\cref{app:two-compression}. The only change is the choice of the two probabilities in~\cref{fig:step} and their common factor $M$. Gaussian sampling and real-valued probability comparisons are ideal operations here. A finite-precision implementation must approximate their joint distribution to negligible statistical distance.

\subsection{Parameters and the Basic Protocol}

\begin{definition}[Signature Parameters]
\label{def:parameters}
Let $d$ be a power of two, let $q$ be an odd prime, and let $\ring\cong\Z[X]/\langle X^d+1\rangle$ and $\ring_a\coloneqq\ring/a\ring$. Fix positive integers $\ell,m,\kappa$ with $\kappa\le d$, positive widths $\sigma,r$, and positive bounds $B_k,B_s$. Write $N\coloneqq d(\ell+m+1)$ and $\vec j\coloneqq(1,0,\ldots,0)^T\in\ring^m$. All norms are Euclidean norms of coefficient vectors. The challenge set $\mathcal C$ is a nonempty subset of binary ring elements of weight at most $\kappa$. The random oracle $H\colon\ring_{2q}^m\times\{0,1\}^*\to\mathcal C$ returns independent uniform challenges. Public parameters, including the admissible functions and $M$, are implicit in the algorithms.
\end{definition}

For the three parameter sets below, $d=256$. For $\kappa\in\{58,80\}$, the challenge set consists of all binary ring elements of weight exactly $\kappa$. For $\kappa=128$, it consists of all binary ring elements of weight less than $128$, together with those of weight $128$ whose constant coefficient is zero. This last set has size $2^{255}$. These are the challenge spaces used in~\cite{GartnerFull,TCHES:CCDGHK24}. The definition specifies the distribution without requiring a particular hash-to-challenge encoding.

\begin{figure}[!htbp]
\centering
\begin{minipage}[t]{.48\linewidth}
\centering
\procedure[linenumbering]{$\mathsf{KeyGen}()$}{
 \mat A_0\leftarrow\Unif(\ring_q^{m\times\ell}) \\
 \textbf{repeat} \\
 \quad\vec s_0\leftarrow D_{\ring^\ell,\sigma} \\
 \quad\vec e\leftarrow D_{\ring^m,\sigma} \\
 \quad\textbf{repeat} \\
 \qquad f_0\leftarrow D_{\ring,\sigma} \\
 \qquad f\coloneqq2f_0+1 \\
 \quad\textbf{until }f\text{ is invertible in }\ring_q \\
 \quad\vec s\coloneqq(f,\vec s_0^T,\vec e^T)^T \\
 \textbf{until }\norm{\vec s}<B_k \\
 \vec b\coloneqq f^{-1}(\mat A_0\vec s_0+\vec e)\bmod q \\
 \mat A_1\coloneqq[q\vec j-2\vec b,\ 2\mat A_0] \\
 \mat A\coloneqq[\mat A_1,\ 2\mat I_m]\bmod2q \\
 \textbf{return }(\mat A,\vec s)
}
\medskip
\procedure[linenumbering]{$\mathsf{Sign}(\mat A,\vec s,\mu)$}{
 \textbf{repeat} \\
 \quad\vec y\leftarrow D_{\ring^{\ell+m+1},r} \\
 \quad\vec w\coloneqq\mat A\vec y\bmod2q \\
 \quad c\coloneqq H(\vec w,\mu) \\
 \quad\vec z\leftarrow\mathsf{RejectSample}_{\vec s}(\vec y,c) \\
 \textbf{until }\vec z\ne\bot\text{ and }\norm{\vec z}\le B_s \\
 \textbf{return }(\vec z,c)
}
\end{minipage}\hfill
\begin{minipage}[t]{.49\linewidth}
\centering
\procedure[linenumbering]{$\mathsf{RejectSample}_{\vec s}(\vec y,c)$}{
 \vec z\coloneqq\vec y \\
 k\coloneqq\norm{c}_1 \\
 \textbf{for }j=0,\ldots,d-1\text{ with }c_j=1 \\
 \quad\vec z\leftarrow R_{X^j\vec s}(\vec z) \\
 \quad\textbf{if }\vec z=\bot \\
 \qquad\textbf{return }\bot \\
 u\leftarrow\Unif([0,1)) \\
 \textbf{if }u\ge M^{k-\kappa} \\
 \quad\textbf{return }\bot \\
 \textbf{return }\vec z
}
\medskip
\procedure[linenumbering]{$\mathsf{Verify}(\mat A,\mu,(\vec z,c))$}{
 \textbf{if }\vec z\notin\ring^{\ell+m+1}\text{ or }c\notin\mathcal C \\
 \quad\textbf{return }0 \\
 \textbf{if }\norm{\vec z}>B_s \\
 \quad\textbf{return }0 \\
 \vec w'\coloneqq\mat A\vec z-qc\vec j\bmod2q \\
 \textbf{return }[H(\vec w',\mu)=c]
}
\end{minipage}
\caption{Key generation, the iterative sampler, and the basic signature protocol.}
\label{fig:signature}
\end{figure}

\cref{fig:signature} specifies key generation, signing, verification, and the entire iterative sampler. Ring elements are lifted coefficientwise when constructing $\mat A$. The public key can be stored as $(\mat A_0,\vec b)$ because the remaining blocks of $\mat A$ are fixed. Every shift $X^j\vec s$ has norm $\norm{\vec s}<B_k$, since multiplication by a monomial is a signed permutation of coefficients. The improved instantiation uses~\cref{eq:functions} with $M=M^*(r/B_k)$ from~\cref{cor:bound}.

For completeness, the original instantiation uses $M=\Mold(r/B_k)$ from~\cref{eq:original}. With
\[
 S_{\vec v}(\vec y)\coloneqq\frac{1}{\rho_r(\vec y)}\sum_{j\ge0}(-1)^j\rho_r(\vec y+j\vec v),
\]
its probabilities are~\cite[Corollary~1]{GartnerFull}
\begin{align*}
 f^{\mathsf{orig}}_{\vec v}(\vec y)&=\frac1M
 \begin{cases}S_{\vec v}(\vec y)&\inner{\vec y}{\vec v}\ge\norm{\vec v}^2,\\1-S_{\vec v}(-\vec y)&\text{otherwise},\end{cases}
 \\
 g^{\mathsf{orig}}_{\vec v}(\vec y)&=\frac1M
 \begin{cases}1-S_{\vec v}(\vec y)&\inner{\vec y}{\vec v}\ge-\norm{\vec v}^2,\\S_{\vec v}(-\vec y)&\text{otherwise}.\end{cases}
\end{align*}
Thus the replacement changes no move, challenge, key, or verification equation.

\begin{lemma}[Basic Correctness]
\label{lem:basic-correctness}
Every signature returned by $\mathsf{Sign}$ in~\cref{fig:signature} is accepted by $\mathsf{Verify}$.
\end{lemma}
\begin{proof}
Key generation gives $\mat A\vec s=q\vec j\bmod2q$. Indeed, its even part vanishes modulo $2q$ by $f\vec b=\mat A_0\vec s_0+\vec e\bmod q$, and $qf\vec j=q\vec j\bmod2q$ because $f=1\bmod2$. A successful iteration produces $\vec z=\vec y+c'\vec s$, where each nonzero coefficient of $c'$ is $1$ or $-1$ and $c'=c\bmod2$. Therefore $\mat A\vec z=\mat A\vec y+qc\vec j\bmod2q$. Verification reconstructs the hashed commitment, and the signing loop enforces the norm bound.
\end{proof}

\subsection{Conditional Unforgeability}

\begin{definition}[Existential Unforgeability]
\label{def:euf-cma}
In the experiment for existential unforgeability under chosen-message attacks (EUF-CMA), an adversary receives the public key, accesses the random oracle and a signing oracle, and wins by producing an accepted signature on a message it never submitted to the signing oracle. A scheme is EUF-CMA secure if every probabilistic polynomial-time adversary wins with negligible probability in the security parameter $\lambda$. A probability is negligible if it decreases faster than every inverse polynomial in $\lambda$. The no-message version (EUF-NMA) provides no signing oracle.
\end{definition}

\begin{definition}[Bounded-Key NTWE Assumption]
\label{def:ntwe}
The bounded-key decisional NTWE assumption asserts that $(\mat A_0,\vec b)$ output by the exact key-generation distribution of~\cref{fig:signature} is computationally indistinguishable from independent uniform elements of $\ring_q^{m\times\ell}\times\ring_q^m$. Indistinguishability means that every probabilistic polynomial-time distinguisher has negligible difference in acceptance probabilities between the two distributions. This definition includes invertibility resampling and the test $\norm{\vec s}<B_k$.
\end{definition}

\begin{definition}[Bimodal Self-Target MSIS Assumption]
\label{def:self-target}
For uniform independent $\mat A_0,\vec b$, construct $\mat A=[q\vec j-2\vec b,2\mat A_0,2\mat I_m]\bmod2q$. Given $\mat A$ and a random oracle $H\colon\ring_{2q}^m\times\{0,1\}^*\to\mathcal C$, the problem at bound $B$ is to find $(\vec z,c,\mu)$ such that
\[
 \vec z\in\ring^{\ell+m+1}\setminus\{\vec0\},\qquad
 \norm{\vec z}\le B,\qquad c\in\mathcal C,\qquad
 H(\mat A\vec z-qc\vec j\bmod2q,\mu)=c.
\]
The assumption asserts that every probabilistic polynomial-time algorithm has negligible success probability. This is the self-target assumption used for ordinary unforgeability in~\cite{GartnerFull}.
\end{definition}

The associated identification protocol uses an independent uniform verifier challenge in place of the hash value. Its accepting transcript can be sampled from the public key by first choosing a norm-truncated Gaussian response and a uniform challenge, then reconstructing the commitment. The full identification protocol and the simulators are given in~\cref{app:security}.

\begin{lemma}[Accepting Simulation and Commitment Entropy]
\label{lem:simulation}
Use any admissible functions with a common factor $M$. Put $\eta\coloneqq\Pr_{\vec z\leftarrow D_{\ring^{\ell+m+1},r}}[\norm{\vec z}\le B_s]$. The identification protocols accept with probability $\eta M^{-\kappa}$, and the simulators of~\cref{fig:simulator} sample their exact transcript distributions conditional on acceptance. For every generated key, either commitment has min-entropy at least
\begin{equation}
\label{eq:entropy-bound}
 \gamma_r\coloneqq d\log_2\bigl(1+1/M^*(r)\bigr).
\end{equation}
Here min-entropy at least $\gamma$ means that no fixed commitment has probability greater than $2^{-\gamma}$ before the challenge is received.
\end{lemma}

The proof is given in~\cref{app:security}. The commitment entropy comes from the coefficient parities, which the commitment determines even after bit-dropping.

\begin{theorem}[Signature Security]
\label{thm:signature-security}
Let the public parameters be functions of $\lambda$. Assume efficient key generation and sampling, admissible move probabilities, $\eta M^{-\kappa}\ge1/\poly(\lambda)$, $\gamma_r=\omega(\log\lambda)$, and negligible $1/|\mathcal C|$. Under the bounded-key NTWE assumption and the bimodal self-target MSIS assumption at bound $B_s$, the algorithms $(\mathsf{KeyGen},\mathsf{Sign},\mathsf{Verify})$ in~\cref{fig:signature} form an EUF-CMA secure signature scheme in the classical random-oracle model. Under the same conditions and self-target assumption at bound $B_v$, the algorithms $(\mathsf{KeyGen},\mathsf{CSign},\mathsf{CVerify})$ in~\cref{fig:signature,fig:compressed} form such a scheme when~\cref{eq:verification-bound} holds. Both statements apply to the improved functions with $M=M^*(r/B_k)$. Negligible total statistical error in implementing the ideal samplers preserves these conclusions.
\end{theorem}

The proof in~\cref{app:security} first replaces the public key by a uniform key, then maps a nonzero forgery to the self-target problem. The accepting-transcript simulator gives the chosen-message reduction. A zero response contributes at most $(q_H+1)/|\mathcal C|$ for $q_H$ classical hash queries.

\cref{thm:signature-security} establishes ordinary unforgeability. Strong unforgeability also excludes new signatures on previously signed messages and requires the additional response-uniqueness analysis at bound $2B_v$ used in~\cite{GartnerFull}.
\FloatBarrier

\FloatBarrier
% !TEX root = main.tex
\section{Four-Move Signatures}
\label{sec:multimodal}

Adding longer odd multiples of the same vector preserves the parity lower bound of~\cref{lem:lower}. Adding a second orthogonal direction changes that bound. The two one-dimensional parity imbalances multiply, and the acceptance functions can attain the resulting bound. We give optimal four-move functions and a complete compressed signature protocol. A monomial rotation supplies the second orthogonal direction without increasing the shift norm. The verification quotient binds half as many challenge bits as the ring degree.

\subsection{An Explicit Four-Move Rule}

\begin{definition}[Four-Move Probabilities]
\label{def:four-moves}
Let $\vec v_1,\vec v_2\in L$ be nonzero orthogonal vectors of equal norm. Abbreviate $b_i(\vec y)\coloneqq b_{\vec v_i}(\vec y)$ and $S_i(\vec y)\coloneqq\widehat S_{\vec v_i}(\vec y)$, using~\cref{def:parity,def:sum}. For a common factor $M$, define
\begin{equation}
\label{eq:four-probabilities}
 \boxed{\displaystyle
 h_{i,-}(\vec y)\coloneqq\frac{S_i(\vec y)}{M(1+b_1(\vec y)b_2(\vec y))},\qquad
 h_{i,+}(\vec y)\coloneqq\frac{S_i(-\vec y)}{M(1+b_1(\vec y)b_2(\vec y))}.}
\end{equation}
The four-move step returns $\vec y\pm\vec v_i$ with probability $h_{i,\pm}(\vec y)$ and rejects with the remaining probability.
\end{definition}

Each numerator is the same weighted alternating sum used for two moves. The only new factor in its denominator is $1+b_1b_2$.

\begin{theorem}[Optimal Four-Move Acceptance]
\label{thm:four-optimal}
Put $\alpha=r/\norm{\vec v_1}$ and define
\begin{equation}
\label{eq:four-rate}
 M_4^*(\alpha)\coloneqq\frac12\left(M^*(\alpha)+\frac1{M^*(\alpha)}\right).
\end{equation}
For $M\ge M_4^*(\alpha)$, the probabilities of~\cref{def:four-moves} are nonnegative, sum to at most one, and satisfy
\begin{equation}
\label{eq:multi-balance}
 \sum_{i=1}^2\bigl(\rho_r(\vec y+\vec v_i)h_{i,-}(\vec y+\vec v_i)
 +\rho_r(\vec y-\vec v_i)h_{i,+}(\vec y-\vec v_i)\bigr)=\frac{\rho_r(\vec y)}M.
\end{equation}
The factor $M_4^*(\alpha)$ is optimal among all acceptance functions using these four moves and preserving $D_{L,r}$ conditional on acceptance with an input-independent factor.
\end{theorem}
\begin{proof}
Nonnegativity follows from~\cref{thm:construction}. Its sum identity gives
\[
 \sum_{i,\pm}h_{i,\pm}(\vec y)=\frac{b_1+b_2}{M(1+b_1b_2)}.
\]
Define $\delta_i=(1-b_i)/(1+b_i)$ and $\delta_0=(M^*(\alpha)-1)/(M^*(\alpha)+1)$. By~\cref{lem:worst}, $|\delta_i|\le\delta_0$. Direct substitution gives
\[
 \frac{b_1+b_2}{1+b_1b_2}=\frac{1-\delta_1\delta_2}{1+\delta_1\delta_2}
 \le\frac{1+\delta_0^2}{1-\delta_0^2}=M_4^*(\alpha),
\]
proving the budget bound.

Orthogonality makes $b_j(\vec y\pm\vec v_i)=b_j(\vec y)$ for $j\ne i$, while~\cref{lem:line} gives $b_i(\vec y\pm\vec v_i)=1/b_i(\vec y)$. Thus the reciprocal denominator, without $M$, at either neighbour in direction $i$ is $b_i/(b_1+b_2)$. The incoming numerator identity of~\cref{thm:construction} is $\rho_r(\vec y)$. The contribution of direction $i$ to~\cref{eq:multi-balance} is consequently $\rho_r(\vec y)b_i/(M(b_1+b_2))$. Summing over $i$ proves the identity.

For optimality, restrict to $\Z\vec v_1+\Z\vec v_2$, which contains the origin and is invariant under every allowed move. Orthogonality makes its Gaussian masses products of one-dimensional masses. The total masses of the two classes determined by the parity of the sum of the two integer coordinates are
\[
 A=\Te(0)^2+\To(0)^2,\qquad B=2\Te(0)\To(0).
\]
Every move exchanges these classes. Summing the incoming identity over the heavier class gives $A/M\le B$, so $M\ge A/B=M_4^*(\alpha)$.
\end{proof}

\begin{corollary}[Four-Move Rejection]
\label{cor:four-rejection}
Define $\delta_0(\alpha)\coloneqq(M^*(\alpha)-1)/(M^*(\alpha)+1)$. With $\delta_0=\delta_0(\alpha)$, the optimal rejection probabilities for two and four moves are respectively
\[
 p_2=\frac{2\delta_0}{1+\delta_0},\qquad
 p_4=\frac{2\delta_0^2}{1+\delta_0^2}.
\]
A common factor $M=M_4^*(r/B_k)$ is valid whenever both orthogonal shifts have the same norm at most $B_k$.
\end{corollary}
\begin{proof}
Substitute the two factors into $1-1/M$. The common-bound statement follows because $M^*(\alpha)$ decreases with $\alpha$ and $x+x^{-1}$ increases for $x\ge1$.
\end{proof}

The per-step comparison in~\cref{tab:four-rates} follows by evaluating~\cref{eq:original,cor:four-rejection}. These are sampler probabilities at fixed $\alpha$, before choosing a challenge space or changing ring parameters.

\begin{table}[!htbp]
\centering
\begin{tabular}{lrrr}
\toprule
$\alpha$ & Original two moves & Optimal two moves & Optimal four moves\\
\midrule
0.8 & $2^{-1.882}$ & $2^{-2.674}$ & $2^{-6.123}$\\
1.0 & $2^{-4.156}$ & $2^{-5.140}$ & $2^{-11.239}$\\
1.2 & $2^{-7.172}$ & $2^{-8.254}$ & $2^{-17.504}$\\
1.5 & $2^{-12.788}$ & $2^{-14.019}$ & $2^{-29.037}$\\
\bottomrule
\end{tabular}
\caption{Per-step rejection at a fixed width-to-shift ratio.}
\label{tab:four-rates}
\end{table}

\subsection{The Ring Construction and Basic Protocol}

\begin{definition}[Four Ring Moves]
\label{def:unit-family}
Use the ring and key parameters of~\cref{def:parameters}, and put
\[
 b\coloneqq d/2,\qquad \omega\coloneqq X^b,\qquad
 U\coloneqq\{1,-1,\omega,-\omega\},\qquad I\coloneqq(\omega-1)\ring.
\]
The allowed moves for $\vec v$ are $\{\vec v,-\vec v,\omega\vec v,-\omega\vec v\}$. Define the folded parity map by
\begin{equation}
\label{eq:fold}
 \mathsf{Fold}(x)_j\coloneqq x_j+x_{j+b}\bmod2
 \qquad(0\le j<b).
\end{equation}
For a generated public key $(\mat A_0,\vec b)$, define
\[
 \mat B_1\coloneqq[-\vec b,\mat A_0],\qquad
 \mat B\coloneqq[\mat B_1,\mat I_m]\bmod q.
\]
Thus $\mat B=2^{-1}\mat A\bmod q$, with $\mat A$ from~\cref{fig:signature}.
\end{definition}

\begin{lemma}[Move Geometry and Key Equations]
\label{lem:unit-geometry}
For every nonzero $\vec v$, the vectors $\vec v$ and $\omega\vec v$ are orthogonal and have equal coefficient norms. Every $u\in U$ equals $1$ modulo $I$. One has $2\ring\subset I$ and $\ring/I\cong\mathbb F_2[X]/\langle X^b-1\rangle$, with quotient map $\mathsf{Fold}$. Every generated key satisfies
\[
 \mat B\vec s=\vec0\pmod q,\qquad f=1\pmod I.
\]
\end{lemma}
\begin{proof}
Multiplication by $\omega$ has an orthogonal coefficient matrix $\mat J$ with $\mat J^2=-\mat I$. Hence $\mat J^T=-\mat J$ and $\inner{\vec v}{\mat J\vec v}=0$. In the quotient by $\omega-1$, one has $\omega=1$ and $\omega^2=-1$, giving $2=0$ and $u=1$ for every $u\in U$. This quotient is exactly the stated binary ring, with coefficient map~\cref{eq:fold}. Finally, $f\vec b=\mat A_0\vec s_0+\vec e\pmod q$ gives $\mat B\vec s=\vec0\pmod q$, while $f=2f_0+1$ and $2\ring\subset I$ give the second key equation.
\end{proof}

\begin{definition}[Four-Move Challenges]
\label{def:folded-challenges}
Choose $1\le\kappa\le b$. Let $\mathcal C_4$ be a nonempty set of binary ring elements supported on coefficients $0,\ldots,b-1$, each of Hamming weight at most $\kappa$. Identify these elements with their $b$-bit vectors. The random oracle
\[
 H_4\colon\ring_q^m\times\{0,1\}^b\times\{0,1\}^*\to\mathcal C_4
\]
returns independent uniform elements. The full binary choice has $\kappa=b$ and $|\mathcal C_4|=2^b$. A weight cap gives $|\mathcal C_4|\le\sum_{j=0}^{\kappa}\binom bj$, while a fixed-weight choice has cardinality $\binom b\kappa$.
\end{definition}

\begin{definition}[Four-Move Iteration]
\label{def:unit-sampler}
Use~\cref{eq:four-probabilities} with common factor $M=M_4^*(r/B_k)$. The step and iterative sampler are given in~\cref{fig:unit-step}. A challenge of actual weight $k$ invokes $k$ steps. The final retention probability $M^{k-\kappa}$ makes the total acceptance probability independent of $k$.
\end{definition}

\begin{figure}[!htbp]
\centering
\begin{minipage}[t]{.52\linewidth}
\centering
\procedure[linenumbering]{$\mathsf{Step}_{\vec v_1,\vec v_2}(\vec y)$}{
 D\coloneqq M(1+b_{\vec v_1}(\vec y)b_{\vec v_2}(\vec y)) \\
 t\leftarrow\Unif([0,1)) \\
 p_{\mathsf{sum}}\coloneqq0 \\
 \textbf{for }(i,e)\in((1,-1),(1,+1),(2,-1),(2,+1)) \\
 \quad p\coloneqq\widehat S_{\vec v_i}(-e\vec y)/D \\
 \quad p_{\mathsf{sum}}\coloneqq p_{\mathsf{sum}}+p \\
 \quad\textbf{if }t<p_{\mathsf{sum}} \\
 \qquad\textbf{return }\vec y+e\vec v_i \\
 \textbf{return }\bot
}
\end{minipage}\hfill
\begin{minipage}[t]{.46\linewidth}
\centering
\procedure[linenumbering]{$\mathsf{Reject}_{4,\vec s}(\vec y,c)$}{
 \vec z\coloneqq\vec y \\
 k\coloneqq\norm{c}_1 \\
 \textbf{for }j=0,\ldots,b-1\text{ with }c_j=1 \\
 \quad\vec v\coloneqq X^j\vec s \\
 \quad\vec z\leftarrow\mathsf{Step}_{\vec v,\omega\vec v}(\vec z) \\
 \quad\textbf{if }\vec z=\bot \\
 \qquad\textbf{return }\bot \\
 t\leftarrow\Unif([0,1)) \\
 \textbf{if }t\ge M^{k-\kappa} \\
 \quad\textbf{return }\bot \\
 \textbf{return }\vec z
}
\end{minipage}
\caption{The four-move step and iterative sampler.}
\label{fig:unit-step}
\end{figure}

The basic signature in~\cref{fig:unit-signature} transmits the full response. The hash binds its image modulo $q$ and the folded parity needed to recover the challenge. The public key is still determined by $(\mat A_0,\vec b)$.

\begin{figure}[!htbp]
\centering
\begin{minipage}[t]{.49\linewidth}
\centering
\procedure[linenumbering]{$\mathsf{KeyGen}_4()$}{
 (\mat A,\vec s)\leftarrow\mathsf{KeyGen}() \\
 \mat B\coloneqq2^{-1}\mat A\bmod q \\
 \textbf{return }(\mat B,\vec s)
}
\par\medskip
\procedure[linenumbering]{$\mathsf{Sign}_4(\mat B,\vec s,\mu)$}{
 \textbf{repeat} \\
 \quad\vec y\leftarrow D_{\ring^{\ell+m+1},r} \\
 \quad\vec w\coloneqq\mat B\vec y\bmod q \\
 \quad\vec u\coloneqq\mathsf{Fold}(y_0) \\
 \quad c\coloneqq H_4(\vec w,\vec u,\mu) \\
 \quad\vec z\leftarrow\mathsf{Reject}_{4,\vec s}(\vec y,c) \\
 \textbf{until }\vec z\ne\bot\text{ and }\norm{\vec z}\le B_s \\
 \textbf{return }(\vec z,c)
}
\end{minipage}\hfill
\begin{minipage}[t]{.49\linewidth}
\centering
\procedure[linenumbering]{$\mathsf{Verify}_4(\mat B,\mu,(\vec z,c))$}{
 \textbf{if }\vec z\notin\ring^{\ell+m+1}\text{ or }c\notin\mathcal C_4 \\
 \quad\textbf{return }0 \\
 \textbf{if }\norm{\vec z}>B_s \\
 \quad\textbf{return }0 \\
 \vec w'\coloneqq\mat B\vec z\bmod q \\
 \vec u'\coloneqq\mathsf{Fold}(z_0)-c\bmod2 \\
 \textbf{return }[H_4(\vec w',\vec u',\mu)=c]
}
\end{minipage}
\caption{Key generation and the basic four-move signature protocol.}
\label{fig:unit-signature}
\end{figure}

\begin{lemma}[Four-Move Correctness]
\label{lem:unit-correctness}
Every signature returned by $\mathsf{Sign}_4$ is accepted by $\mathsf{Verify}_4$.
\end{lemma}
\begin{proof}
A successful sampler produces $\vec z=\vec y+c'\vec s$, where each nonzero challenge coefficient is replaced by an element of $U$. Thus $c'=c\pmod I$. By~\cref{lem:unit-geometry},
\[
 \mat B\vec z=\mat B\vec y\pmod q,\qquad
 \mathsf{Fold}(z_0)-c=\mathsf{Fold}(y_0)\pmod2.
\]
These equations recover both hash inputs. The signing loop enforces the norm test.
\end{proof}

The compressed protocol and its worst-case reconstruction bound are given in~\cref{app:four-compression}. Bit-dropping leaves the first response ring element unchanged, so it preserves the folded parity used in verification.

\subsection{Conditional Unforgeability}

\begin{lemma}[Simulation and Commitment Entropy]
\label{lem:four-simulation}
Let $\eta$ be the probability that a response sampled from $D_{\ring^{\ell+m+1},r}$ has norm at most $B_s$. Both associated identification protocols accept with probability $\eta M^{-\kappa}$. Their exact accepting transcript distributions are sampled by~\cref{fig:four-simulator}. Before the challenge, either commitment has min-entropy at least
\begin{equation}
\label{eq:fold-entropy}
 \gamma_4\coloneqq b\log_2\frac{2}{1+\delta_0(r)^2}.
\end{equation}
\end{lemma}

The proof and full simulators are given in~\cref{app:security}. Folding pairs independent coefficient parities, so the bias of a folded bit is the square of a single coefficient's parity bias.

\begin{definition}[Folded Self-Target Assumption]
\label{def:quotient-self-target}
Sample independent uniform $\mat A_0,\vec b$, and construct $\mat B$ as in~\cref{def:unit-family}. Given $\mat B$ and the random oracle $H_4$, the folded self-target problem at bound $B$ asks for
\[
 \vec z\in\ring^{\ell+m+1}\setminus\{\vec0\},\qquad
 \norm{\vec z}\le B,\qquad c\in\mathcal C_4,\qquad\mu\in\{0,1\}^*
\]
such that
\[
 H_4(\mat B\vec z\bmod q,\mathsf{Fold}(z_0)-c\bmod2,\mu)=c.
\]
The assumption asserts negligible success for every probabilistic polynomial-time algorithm.
\end{definition}

\begin{theorem}[Four-Move Signature Security]
\label{thm:unit-security}
Assume efficient key generation and sampling, the bounded-key NTWE assumption, $\eta M^{-\kappa}\ge1/\poly(\lambda)$, $\gamma_4=\omega(\log\lambda)$, and negligible $1/|\mathcal C_4|$. Under the folded self-target assumption at bound $B_s$, the algorithms $(\mathsf{KeyGen}_4,\mathsf{Sign}_4,\mathsf{Verify}_4)$ form an EUF-CMA secure signature scheme in the classical random-oracle model. Under the same conditions and the folded self-target assumption at bound $B_v$, the algorithms $(\mathsf{KeyGen}_4,\mathsf{CSign}_4,\mathsf{CVerify}_4)$ form such a scheme when~\cref{eq:four-verification-bound} holds. These conclusions use $M=M_4^*(r/B_k)$ and remain valid with negligible total statistical error in implementing the ideal samplers. Replacing the four-move step by either admissible two-move rule, with its common factor $M$, gives the same conclusions.
\end{theorem}

The proof is given in~\cref{app:security}. The folded self-target assumption concerns a different relation from~\cref{def:self-target}. The theorem states ordinary unforgeability under this assumption. The concrete comparison in~\cref{sec:results} uses this same folded relation for all three acceptance rules within each parameter set.

\begin{remark}[Additional Unit Moves]
\label{rem:more-moves}
For a power of two $a\mid d$, the $2a$ moves $\{\pm X^{jd/a}\vec v:0\le j<a\}$ split into $a/2$ orthogonal pairs of directions. Uniformly choosing a pair and applying~\cref{eq:four-probabilities} has the same factor $M_4^*$, since averaging the incoming identities preserves that factor. The quotient identifying all units with one has only $d/a$ challenge bits, because setting $X^{d/a}=1$ gives $\mathbb F_2[X]/\langle X^{d/a}-1\rangle$. Thus this extension gives no further rejection gain and reduces challenge capacity. Mutual orthogonality of all directions does not follow from equal norms. For eight moves, put $\xi=X^{d/4}$ and $v=1+\xi$. Then $\norm{v}^2=2$ but $\inner{v}{\xi v}=1$. This does not exclude an improvement from a different acceptance rule.
\end{remark}

\FloatBarrier

\FloatBarrier
% !TEX root = main.tex
\section{Signature Sizes and Repetition Counts}
\label{sec:results}

We compare acceptance rules in two settings. The first uses the original rejecting parameter sets of~\cite{GartnerFull} without changing the signature relation. The second fixes three parameter sets for the folded relation and compares all three rules within that relation. The byte counts use the entropy model of~\cite[Equation~(8) and Appendix~A.2]{GartnerFull}, including challenge bytes and padding. Supporting calculations are given in~\cref{sec:comparison,app:four-comparison}.

\subsection{The Original Two-Move Parameter Sets}

\cref{tab:summary-two} gives the parameters and both ways to use optimal two-move acceptance. The size comparison lowers the Gaussian width while preserving the original iterative-sampler repetition factor. The attempt comparison retains the original width, so its signature-size estimate is unchanged. All comparisons use the same keys, challenge distributions, and compression parameters within each set.

\begin{table}[!htbp]
\centering
\begin{tabular}{lrrr}
\toprule
Parameter or estimate & I & II & III\\
\midrule
Ring degree $d$ & 256 & 256 & 256\\
Modulus $q$ & 12289 & 50177 & 50177\\
Module parameters $(\ell,m)$ & $(1,2)$ & $(3,2)$ & $(4,3)$\\
SIS coefficient matrix & $512\times1024$ & $512\times1536$ & $768\times2048$\\
Secret width $\sigma$ & 2.60 & 1.00 & 1.50\\
Secret bound $B_k$ & 110.07 & 48.00 & 79.60\\
Challenge HW bound $\kappa$ & 58 & 80 & 128\\
Compression parameter $\tau$ & 256 & 128 & 128\\
Original response width & 128 & 55 & 95\\
Matched optimal width & 120.6834 & 51.7805 & 89.7892\\
\midrule
Original signature bytes & 775 & 1184 & 1694\\
Optimal bytes at matched rejection & 764 & 1167 & 1674\\
Original sampler attempts & 1.850 & 2.790 & 2.612\\
Optimal attempts at original width & 1.341 & 1.634 & 1.574\\
\bottomrule
\end{tabular}
\caption{Two-move parameters and improvements at the parameter sets of~\cite{GartnerFull}.}
\label{tab:summary-two}
\end{table}

Sets I and II use fixed binary Hamming weights $58$ and $80$. Set III uses weights at most $128$ and the half-space selection of~\cref{def:parameters}, giving $2^{255}$ challenges. None uses full binary challenges. The detailed reproduction in~\cref{tab:parameters} matches all three original signature sizes and all nine rejecting-mode BKZ and Core-SVP estimate pairs. It retains the listed set-III verification bound $5300$ for that reproduction. Applying the displayed integer rounding bound requires $5301$, as accounted for separately in~\cref{tab:bounds}.

\subsection{Four Moves at Fixed Rejection}

\cref{tab:four-parameters} gives the three parameter sets for the folded relation. All use $q=50177$ and rounding step $\Delta=64$. Sets A and C use fixed Hamming weights $58$ on $256$ positions and $82$ on $512$ positions. Set B uses all $256$-bit binary challenges, with maximum weight $256$. Its mean weight $128$ is not the repetition exponent. The key distribution, folded verification relation, challenge distribution, and encoding model are fixed within each comparison. In the table, $h=\log_2|\mathcal C_4|$ is the challenge entropy, and $t$ allocates $2^{-t}$ to the challenge-search query term at $Q=2^{64}$ classical queries, as formalised in~\cref{def:query-budget}.
\begin{table}[!htbp]
\centering
\begin{tabular}{lrrr}
\toprule
Parameter & A & B & C\\
\midrule
Ring degree $d$ & 512 & 512 & 1024\\
Module parameters $(\ell,m)$ & $(2,2)$ & $(3,2)$ & $(1,2)$\\
Module matrix $\mat B$ & $2\times5$ & $2\times6$ & $2\times4$\\
SIS coefficient matrix & $1024\times2560$ & $1024\times3072$ & $2048\times4096$\\
NTWE proxy dimension & 1535 & 2047 & 2047\\
Secret width $\sigma$ & 0.90 & 1.00 & 1.00\\
Secret bound $B_k$ & 58 & 68 & 85\\
Challenge positions $b$ & 256 & 256 & 512\\
Weight rule & Fixed & Full binary & Fixed\\
Maximum weight $\kappa$ & 58 & 256 & 82\\
Challenge entropy $h$ & 193.553 & 256 & 320.578\\
Query target exponent $t$ & 128 & 192 & 256\\
Query contribution $Q2^{-h}$ & $2^{-129.553}$ & $2^{-192}$ & $2^{-256.578}$\\
Challenge bytes & 32 & 32 & 41\\
Key bytes & 2080 & 2080 & 4128\\
\bottomrule
\end{tabular}
\caption{Comparison parameters with $q=50177$, $\Delta=64$, and $Q=2^{64}$.}
\label{tab:four-parameters}
\end{table}

\cref{tab:four-sizes} compares the two-move rule of~\cite{GartnerFull}, optimal two moves, and optimal four moves. Every entry has two expected sampler attempts before norm rejection. Of the total savings of $162$, $195$, and $259$ bytes, optimal two moves account for $28$, $28$, and $42$ bytes. Adding the second direction saves the remaining $134$, $167$, and $217$ bytes.
\begin{table}[!htbp]
\centering
\begin{tabular}{lrrrrr}
\toprule
& \multicolumn{3}{c}{Estimated signature bytes} & \multicolumn{2}{c}{Saved}\\
Set & Rule of~\cite{GartnerFull} & Optimal two & Four & Bytes & Fraction\\
\midrule
A & 1892 & 1864 & 1730 & 162 & $8.56\%$\\
B & 2557 & 2529 & 2362 & 195 & $7.63\%$\\
C & 2991 & 2949 & 2732 & 259 & $8.66\%$\\
\bottomrule
\end{tabular}
\caption{Signature sizes at fixed sampler rejection across three parameter sets.}
\label{tab:four-sizes}
\end{table}

\cref{tab:four-bounds} gives every response width and norm bound used in the size comparison. The signing cutoff is $B_s=\lceil1.02r\sqrt N\rceil$, and the verification cutoff is $B_v=\lceil B_s+\sqrt{dm}(\Delta/2+1)\rceil$, where $N=d(\ell+m+1)$. The second term bounds the compression error in the worst case. All three rules use the same folded relation, so~\cref{thm:unit-security} applies under the folded self-target assumption at the corresponding verification bound.
\begin{table}[!htbp]
\centering
\begin{tabular}{llrrrr}
\toprule
Set & Acceptance rule & $r$ & $B_s$ & $B_v$ & $2B_v$\\
\midrule
A & Rule of~\cite{GartnerFull} & 66.82808338 & 3449 & 4505 & 9010\\
  & Optimal two & 62.95106391 & 3249 & 4305 & 8610\\
  & Four & 47.09214965 & 2431 & 3487 & 6974\\
\midrule
B & Rule of~\cite{GartnerFull} & 86.98084173 & 4918 & 5974 & 11948\\
  & Optimal two & 82.69429158 & 4676 & 5732 & 11464\\
  & Four & 61.18756830 & 3460 & 4516 & 9032\\
\midrule
C & Rule of~\cite{GartnerFull} & 100.55658966 & 6565 & 8059 & 16118\\
  & Optimal two & 94.96376526 & 6200 & 7694 & 15388\\
  & Four & 70.82704566 & 4624 & 6118 & 12236\\
\bottomrule
\end{tabular}
\caption{Response widths and Euclidean norm bounds at fixed sampler rejection.}
\label{tab:four-bounds}
\end{table}

At the original-rule widths instead, all three rules have the same estimated size. Four moves reduce expected sampler attempts from $2$ to $1.000947$, $1.000199$, and $1.000658$. The corresponding optimal two-move counts are $1.392827$, $1.375682$, and $1.388798$. These are repetition counts, not running times, since four moves require evaluating an additional direction. The complete fixed-width comparison is given in~\cref{tab:four-fixed-width}.

\subsection{Challenge Capacity and Norm Accounting}

The query budget in~\cref{tab:four-parameters} is an illustrative maximum of $Q=2^{64}$ classical hash queries, distinct from the modulus $q$. With challenge entropy $h=\log_2|\mathcal C_4|$, the query contribution is $Q2^{-h}$. Thus the listed target exponents allocate only the challenge-search term. They are not complete signature-security levels. The full bound includes the folded self-target and bounded-key NTWE advantages and the signing-simulation terms of~\cref{lem:fs-transfer}. The exact challenge-search bound, including an unqueried final forgery, is given in~\cref{prop:challenge-search}.

The accepted Gaussian response law, rather than an average-case estimate of challenge multiplication, supplies the width used in the entropy calculation. Every signature also satisfies a hard norm test. The raw shift bound is $\norm{c'\vec s}\le\kappa B_k$, but this is not added to the Gaussian radius in the size model. The fold acts only on commitment parity. The separate Gaussian norm-tail bound is given in~\cref{lem:multi-norm-tail}. Two expected sampler attempts do not include norm or encoding rejection.

The mathematical security statements use exact Gaussian sampling and exact acceptance probabilities. A finite implementation must control their total statistical error and specify its byte encoding and overflow rejection. The unstructured lattice estimates in~\cref{tab:four-lattice} do not establish the folded self-target assumption.

\FloatBarrier

\bibliography{../optimal-gartner/cryptobib/abbrev3,../optimal-gartner/cryptobib/crypto,biblio}
\clearpage
\appendix
% !TEX root = main.tex
\section{Bit-Dropping and Compressed Protocols}
\label{app:compression}

The protocols replace the last $m$ response ring elements by rounding hints. The reconstruction error is bounded coefficientwise in the worst case. We give both variants, including their full verification checks.

\subsection{Two-Move Compression}
\label{app:two-compression}

\begin{definition}[Rounding and Hints]
\label{def:rounding}
Choose an integer $\tau\ge4$ divisible by $4$ and dividing $2(q-1)$. Write $x\bmod a$ for the representative in $\{0,\ldots,a-1\}$ and $x\bmod^{\pm}a$ for the representative in $[-a/2,a/2)$. Define coefficientwise
\begin{align*}
 \mathsf{HB}(x)&\coloneqq\tau\left\lfloor\frac{x\bmod2q}{\tau}+\frac12\right\rfloor\bmod2(q-1),\\
 \mathsf{LB}(x)&\coloneqq x-\mathsf{HB}(x)\bmod^{\pm}2q,
 &\mathsf{LSB}(x)&\coloneqq x\bmod2.
\end{align*}
Let $\mathcal G_\tau\coloneqq\{0,\tau,\ldots,2(q-1)-\tau\}$, and let $\mathcal H_\tau$ be the set of vectors of $m$ ring elements whose coefficients lie in $\mathcal G_\tau$. The compressed random oracle is $H_c\colon\mathcal H_\tau\times\ring_2^m\times\{0,1\}^*\to\mathcal C$. Fix an integer verification bound
\begin{equation}
\label{eq:verification-bound}
 B_v\ge B_s+\sqrt{dm}\,(\tau/4+1).
\end{equation}
\end{definition}

\begin{figure}[!htbp]
\centering
\procedure[linenumbering]{$\mathsf{CSign}(\mat A,\vec s,\mu)$}{
 \textbf{repeat} \\
 \quad\vec y\leftarrow D_{\ring^{\ell+m+1},r} \\
 \quad\vec w\coloneqq\mat A\vec y\bmod2q \\
 \quad\vec u\coloneqq\mathsf{LSB}(\vec w) \\
 \quad c\coloneqq H_c(\mathsf{HB}(\vec w),\vec u,\mu) \\
 \quad\vec z\leftarrow\mathsf{RejectSample}_{\vec s}(\vec y,c) \\
 \textbf{until }\vec z\ne\bot\text{ and }\norm{\vec z}\le B_s \\
 (\vec z_1,\vec z_2)\coloneqq\mathsf{Split}(\vec z) \\
 \widetilde{\vec w}\coloneqq\vec w-2\vec z_2\bmod2q \\
 \vec h\coloneqq\mathsf{HB}(\vec w)-\mathsf{HB}(\widetilde{\vec w})\bmod2(q-1) \\
 \textbf{return }(\vec z_1,\vec h,c)
}
\par\medskip
\procedure[linenumbering]{$\mathsf{CVerify}(\mat A,\mu,(\vec z_1,\vec h,c))$}{
 \textbf{if }\vec z_1\notin\ring^{\ell+1}\text{ or }\vec h\notin\mathcal H_\tau\text{ or }c\notin\mathcal C \\
 \quad\textbf{return }0 \\
 \widetilde{\vec w}'\coloneqq\mat A_1\vec z_1-qc\vec j\bmod2q \\
 \vec{w}_{h}'\coloneqq\mathsf{HB}(\widetilde{\vec w}')+\vec h\bmod2(q-1) \\
 z_0\coloneqq\text{first ring element of }\vec z_1 \\
 \vec u'\coloneqq\mathsf{LSB}(z_0-c)\vec j \\
 \vec z_2'\coloneqq(\vec{w}_{h}'-\widetilde{\vec w}'+\vec u')/2\bmod^{\pm}q \\
 \vec z'\coloneqq(\vec z_1^T,(\vec z_2')^T)^T \\
 \textbf{if }\norm{\vec z'}>B_v \\
 \quad\textbf{return }0 \\
 \textbf{return }[H_c(\vec{w}_{h}',\vec u',\mu)=c]
}
\caption{Compressed signing and verification with the same key generation and sampler.}
\label{fig:compressed}
\end{figure}

In~\cref{fig:compressed}, $\mathsf{Split}$ separates the first $\ell+1$ ring elements from the last $m$, and $\mat A_1$ is the corresponding block of $\mat A$. The division by $2$ is an integer division before reduction modulo $q$. It is always defined because $\widetilde{\vec w}'\bmod2=\vec u'$ and $\vec{w}_{h}'$ has even coefficients. The hint is represented canonically, not as an arbitrary congruent integer vector. These algorithms specify mathematical signatures. A byte encoding must additionally specify an injective encoding of $\vec z_1$, $\vec h/\tau$, and $c$, and reject malformed encodings. The entropy estimates below are not a specification of an entropy coder.

\begin{lemma}[Compressed Correctness]
\label{lem:compressed-correctness}
Under~\cref{eq:verification-bound}, every signature returned by $\mathsf{CSign}$ in~\cref{fig:compressed} is accepted by $\mathsf{CVerify}$.
\end{lemma}
\begin{proof}
By~\cref{lem:basic-correctness}, $\widetilde{\vec w}'=\vec w-2\vec z_2\bmod2q$, so the hint gives $\vec{w}_{h}'=\mathsf{HB}(\vec w)$. Also $\mat A\bmod2=[\vec j,0,\ldots,0]$ and $z_0-c=y_0\bmod2$, giving $\vec u'=\mathsf{LSB}(\vec w)$. Thus the hash inputs agree.

Put $\vec l=\mathsf{LB}(\vec w)$ and $\vec u=\mathsf{LSB}(\vec w)$. Ordinary rounding gives coefficients in $[-\tau/2,\tau/2-1]$. When the rounded value wraps at $2(q-1)$, centring modulo $2q$ introduces an additional $-2$. Consequently every coefficient of $\vec l$ lies in $[-\tau/2-2,\tau/2-1]$. Since $\vec l\bmod2=\vec u$ and $4$ divides $\tau$, every coefficient of $(\vec u-\vec l)/2$ has magnitude at most $\tau/4+1$. Reconstruction gives
\[
 \vec z_2'=\vec z_2+(\vec u-\vec l)/2\bmod^{\pm}q.
\]
Centred reduction cannot increase the magnitude of an integer coefficient. The triangle inequality now gives $\norm{\vec z'}\le\norm{\vec z}+\sqrt{dm}(\tau/4+1)\le B_v$.
\end{proof}

\FloatBarrier
\subsection{Four-Move Compression}
\label{app:four-compression}

The last $m$ ring elements can again be replaced by rounding hints. Using $\mat B$ makes their coefficient in the verification equation equal to one. This removes the parity correction needed when dividing by two in~\cref{fig:compressed}.

\begin{definition}[Four-Move Rounding]
\label{def:four-rounding}
Choose an even integer $\Delta\ge2$ dividing $q-1$. Define coefficientwise
\[
 \mathsf{Q}_\Delta(x)\coloneqq
 \Delta\left\lfloor\frac{x\bmod q}{\Delta}+\frac12\right\rfloor\bmod(q-1),
 \qquad
 \mathsf{L}_\Delta(x)\coloneqq x-\mathsf{Q}_\Delta(x)\bmod^{\pm}q.
\]
Let $\mathcal H_\Delta$ consist of vectors of $m$ ring elements whose coefficients belong to $\{0,\Delta,\ldots,q-1-\Delta\}$. Fix an integer bound
\begin{equation}
\label{eq:four-verification-bound}
 B_v\ge B_s+\sqrt{dm}\,(\Delta/2+1).
\end{equation}
\end{definition}

\cref{fig:four-compressed} gives the compressed protocol. Its signature is $(\vec z_1,\vec h,c)$, where $\vec z_1$ contains the first $\ell+1$ ring elements. The folded parity is reconstructed from $z_0$ and $c$, so it costs no additional signature bits. The hint is a canonical element of $\mathcal H_\Delta$.

\begin{figure}[!htbp]
\centering
\procedure[linenumbering]{$\mathsf{CSign}_4(\mat B,\vec s,\mu)$}{
 \textbf{repeat} \\
 \quad\vec y\leftarrow D_{\ring^{\ell+m+1},r} \\
 \quad\vec w\coloneqq\mat B\vec y\bmod q \\
 \quad\vec u\coloneqq\mathsf{Fold}(y_0) \\
 \quad c\coloneqq H_4(\mathsf{Q}_\Delta(\vec w),\vec u,\mu) \\
 \quad\vec z\leftarrow\mathsf{Reject}_{4,\vec s}(\vec y,c) \\
 \textbf{until }\vec z\ne\bot\text{ and }\norm{\vec z}\le B_s \\
 (\vec z_1,\vec z_2)\coloneqq\mathsf{Split}(\vec z) \\
 \widetilde{\vec w}\coloneqq\vec w-\vec z_2\bmod q \\
 \vec h\coloneqq\mathsf{Q}_\Delta(\vec w)-\mathsf{Q}_\Delta(\widetilde{\vec w})\bmod(q-1) \\
 \textbf{return }(\vec z_1,\vec h,c)
}
\par\medskip
\procedure[linenumbering]{$\mathsf{CVerify}_4(\mat B,\mu,(\vec z_1,\vec h,c))$}{
 \textbf{if }\vec z_1\notin\ring^{\ell+1}\text{ or }\vec h\notin\mathcal H_\Delta\text{ or }c\notin\mathcal C_4 \\
 \quad\textbf{return }0 \\
 \widetilde{\vec w}'\coloneqq\mat B_1\vec z_1\bmod q \\
 \vec w_h'\coloneqq\mathsf{Q}_\Delta(\widetilde{\vec w}')+\vec h\bmod(q-1) \\
 \vec z_2'\coloneqq\vec w_h'-\widetilde{\vec w}'\bmod^{\pm}q \\
 \vec z'\coloneqq(\vec z_1^T,(\vec z_2')^T)^T \\
 \textbf{if }\norm{\vec z'}>B_v \\
 \quad\textbf{return }0 \\
 z_0\coloneqq\text{first ring element of }\vec z_1 \\
 \vec u'\coloneqq\mathsf{Fold}(z_0)-c\bmod2 \\
 \textbf{return }[H_4(\vec w_h',\vec u',\mu)=c]
}
\caption{Compressed signing and verification for four moves.}
\label{fig:four-compressed}
\end{figure}

\begin{lemma}[Compressed Four-Move Correctness]
\label{lem:four-compressed}
Under~\cref{eq:four-verification-bound}, every signature returned by $\mathsf{CSign}_4$ is accepted by $\mathsf{CVerify}_4$. Its reconstructed response satisfies
\[
 \mat B\vec z'=\vec w_h'\pmod q,\qquad
 \norm{\vec z'}\le B_v.
\]
\end{lemma}
\begin{proof}
The key equation gives $\widetilde{\vec w}'=\vec w-\vec z_2\pmod q$, so the hint recovers $\vec w_h'=\mathsf{Q}_\Delta(\vec w)$. As in~\cref{lem:unit-correctness}, the folded parity is recovered exactly. Thus both hash inputs agree.

For a scalar residue, ordinary rounding has error in $[-\Delta/2,\Delta/2-1]$. A wrap at $q-1$ gives a congruent error smaller by one modulo $q$. Thus every coefficient of $\vec l=\mathsf{L}_\Delta(\vec w)$ is the centred reduction of an integer in $[-\Delta/2-1,\Delta/2-1]$, and has magnitude at most $\Delta/2+1$. Reconstruction gives
\[
 \vec z_2'=\vec z_2-\vec l\bmod^{\pm}q.
\]
Centred reduction cannot increase coefficient magnitude. Consequently
\[
 \norm{\vec z'}\le\norm{\vec z}+\sqrt{dm}\,(\Delta/2+1)\le B_v.
\]
Finally, $\mat B=[\mat B_1,\mat I_m]$ gives the claimed reconstructed commitment equation directly.
\end{proof}

\FloatBarrier

% !TEX root = main.tex
\section{Simulation and Security Proofs}
\label{app:security}

\begin{definition}[Associated Identification Protocol]
\label{def:identification}
The associated three-message identification protocol replaces the hash call in a single signing attempt by an independent uniform challenge from an honest verifier. The prover first sends $\vec w$ in the basic variant, or $(\mathsf{HB}(\vec w),\mathsf{LSB}(\vec w))$ in the compressed variant. After receiving $c$, it runs the sampler, aborts on failure or excessive response norm, and sends $\vec z$ or $(\vec z_1,\vec h)$. The verifier checks the corresponding reconstruction equation and norm bound. Its transcript is the commitment, challenge, and response. The Fiat-Shamir transformation replaces this challenge with the random-oracle value on the commitment and message, giving~\cref{fig:signature,fig:compressed}.
\end{definition}

\subsection{Two-Move Transcript Simulation}

\begin{figure}[t]
\centering
\begin{minipage}[t]{.44\linewidth}
\centering
\procedure[linenumbering]{$\mathsf{Sim}(\mat A)$}{
 \textbf{repeat} \\
 \quad\vec z\leftarrow D_{\ring^{\ell+m+1},r} \\
 \textbf{until }\norm{\vec z}\le B_s \\
 c\leftarrow\Unif(\mathcal C) \\
 \vec w\coloneqq\mat A\vec z-qc\vec j\bmod2q \\
 \textbf{return }(\vec w,c,\vec z)
}
\end{minipage}\hfill
\begin{minipage}[t]{.54\linewidth}
\centering
\procedure[linenumbering]{$\mathsf{CSim}(\mat A)$}{
 (\vec w,c,\vec z)\leftarrow\mathsf{Sim}(\mat A) \\
 (\vec z_1,\vec z_2)\coloneqq\mathsf{Split}(\vec z) \\
 \vec{w}_{h}\coloneqq\mathsf{HB}(\vec w) \\
 \vec u\coloneqq\mathsf{LSB}(\vec w) \\
 \widetilde{\vec w}\coloneqq\vec w-2\vec z_2\bmod2q \\
 \vec h\coloneqq\vec{w}_{h}-\mathsf{HB}(\widetilde{\vec w})\bmod2(q-1) \\
 \textbf{return }((\vec{w}_{h},\vec u),c,(\vec z_1,\vec h))
}
\end{minipage}
\caption{Public-key simulators for accepting identification transcripts.}
\label{fig:simulator}
\end{figure}

\subsubsection{Two-Move Transcript Distribution}

\begin{proof}[Proof of~\cref{lem:simulation}]
Fix $c$. Its monomial shifts are fixed independently of $\vec y$, so~\cref{cor:iteration} gives the unnormalised response law $M^{-\kappa}D_{\ring^{\ell+m+1},r}$, irrespective of $c$. The norm test multiplies acceptance by $\eta$, also independently of $c$. Conditional on acceptance, the challenge is uniform and independent of the norm-truncated Gaussian response. By~\cref{lem:basic-correctness}, the commitment is determined by $(\vec z,c)$ as $\vec w=\mat A\vec z-qc\vec j\bmod2q$. This proves exact simulation in the basic case. The compressed transcript is the deterministic transformation implemented by $\mathsf{CSim}$, proving the compressed case.

Both commitments determine $\mathsf{LSB}(\vec w)=\mathsf{LSB}(y_0)\vec j$. The $d$ coefficients of $y_0$ are independent samples from $D_{\Z,r}$. Their larger parity probability is the even probability $M^*(r)/(1+M^*(r))$, by the definition of the centred parity masses. Every fixed parity vector therefore has probability at most $(1+1/M^*(r))^{-d}$. Specifying the whole commitment cannot increase this probability, proving~\cref{eq:entropy-bound}.
\end{proof}

\subsection{Fiat-Shamir Transfer and Two-Move Security}

\begin{lemma}[Fiat-Shamir Security Transfer, {\cite[Theorem~3]{GartnerFull}}]
\label{lem:fs-transfer}
Consider a three-message identification protocol with commitment min-entropy at least $\gamma$, abort probability $\beta<1$, and an efficient public-key simulator for the exact distribution of accepting honest-verifier transcripts. For every classical EUF-CMA adversary making $q_S$ signing queries and $q_H$ random-oracle queries against its Fiat-Shamir signature scheme, there is an EUF-NMA adversary whose success probability $\epsilon_{\mathsf{NMA}}$ satisfies
\[
 \epsilon_{\mathsf{CMA}}\le\epsilon_{\mathsf{NMA}}
 +\frac{2q_S2^{-\gamma/2}}{1-\beta}\sqrt{q_H+1+\frac{q_S}{1-\beta}}
 +2^{1-\gamma/2}(q_H+1)\sqrt{\frac{q_S}{1-\beta}}.
\]
Here $\epsilon_{\mathsf{CMA}}$ is the chosen-message adversary's success probability. The no-message adversary's running time is that of the chosen-message adversary plus the cost of $q_S$ simulator calls, up to polynomial overhead.
\end{lemma}

\subsubsection{Two-Move Unforgeability}

\begin{proof}[Proof of~\cref{thm:signature-security}]
Correctness is given by~\cref{lem:basic-correctness,lem:compressed-correctness}. First consider an adversary with no signing queries. Replacing the generated public key by uniform $(\mat A_0,\vec b)$ changes its success probability negligibly under~\cref{def:ntwe}. For the basic scheme, every forgery with nonzero response is then a solution of~\cref{def:self-target} at bound $B_s$.

For the compressed scheme, provide the adversary with $H_c(\vec{w}_{h},\vec u,\mu)=H(\vec{w}_{h}+\vec u\bmod2q,\mu)$. This is a random oracle on its stated domain because the map $(\vec{w}_{h},\vec u)\mapsto\vec{w}_{h}+\vec u\bmod2q$ is injective. Indeed, each coefficient of $\vec{w}_{h}$ is a canonical multiple of $\tau$, and adding a bit neither overlaps adjacent multiples nor wraps modulo $2q$. From any accepted forgery, reconstruct $\vec z'$ as in~\cref{fig:compressed}. Its norm is at most $B_v$, and
\[
 \mat A\vec z'-qc\vec j=\vec{w}_{h}'+\vec u'\pmod{2q}.
\]
Thus a nonzero $\vec z'$ solves the same self-target problem at bound $B_v$.

A zero response does not solve the stated nonzero-vector problem. However, it requires $H(-qc\vec j,\mu)=c$. The map $c\mapsto-qc\vec j\bmod2q$ is injective on binary challenges. Each classical oracle query can therefore succeed for at most one predetermined challenge, with probability $1/|\mathcal C|$. Including an unqueried final forgery, the zero-response contribution is at most $(q_H+1)/|\mathcal C|$ for $q_H$ queries. The compressed embedding has the same property. This completes the no-message security argument.

Apply~\cref{lem:fs-transfer}. Its hypotheses hold by~\cref{lem:simulation} with abort probability $\beta=1-\eta M^{-\kappa}$ and entropy bound $\gamma_r$. The simulator has expected polynomial running time because $\eta$ is inverse polynomial. For polynomial query counts, both additive terms are negligible because $\gamma_r=\omega(\log\lambda)$ and $(1-\beta)^{-1}$ is polynomially bounded. Hence the chosen-message advantage is negligible. Finally,~\cref{thm:construction,cor:bound} provide admissibility for the improved functions. A coupling with total negligible implementation error changes the adversary's success probability negligibly.
\end{proof}

\subsection{Four-Move Transcript Simulation}

\cref{fig:four-simulator} specifies the accepting-transcript simulators. The compressed simulator applies the same rounding and hint map as the signer.

\begin{figure}[!htbp]
\centering
\begin{minipage}[t]{.44\linewidth}
\centering
\procedure[linenumbering]{$\mathsf{Sim}_4(\mat B)$}{
 \textbf{repeat} \\
 \quad\vec z\leftarrow D_{\ring^{\ell+m+1},r} \\
 \textbf{until }\norm{\vec z}\le B_s \\
 c\leftarrow\Unif(\mathcal C_4) \\
 \vec w\coloneqq\mat B\vec z\bmod q \\
 \vec u\coloneqq\mathsf{Fold}(z_0)-c\bmod2 \\
 \textbf{return }((\vec w,\vec u),c,\vec z)
}
\end{minipage}\hfill
\begin{minipage}[t]{.54\linewidth}
\centering
\procedure[linenumbering]{$\mathsf{CSim}_4(\mat B)$}{
 ((\vec w,\vec u),c,\vec z)\leftarrow\mathsf{Sim}_4(\mat B) \\
 (\vec z_1,\vec z_2)\coloneqq\mathsf{Split}(\vec z) \\
 \vec w_h\coloneqq\mathsf{Q}_\Delta(\vec w) \\
 \widetilde{\vec w}\coloneqq\vec w-\vec z_2\bmod q \\
 \vec h\coloneqq\vec w_h-\mathsf{Q}_\Delta(\widetilde{\vec w})\bmod(q-1) \\
 \textbf{return }((\vec w_h,\vec u),c,(\vec z_1,\vec h))
}
\end{minipage}
\caption{Accepting-transcript simulators for the four-move protocols.}
\label{fig:four-simulator}
\end{figure}

\subsubsection{Four-Move Transcript Distribution}

\begin{proof}[Proof of~\cref{lem:four-simulation}]
For each fixed challenge,~\cref{thm:four-optimal} gives the incoming Gaussian identity at each step. The final retention test gives unnormalised response law $M^{-\kappa}D_{\ring^{\ell+m+1},r}$, independently of the challenge weight. The norm test multiplies acceptance by $\eta$. Conditional on acceptance, $\vec z$ is a norm-truncated Gaussian independent of the uniform challenge. The key equations determine the basic commitment from $(\vec z,c)$, proving exact simulation. The compressed transcript is the deterministic transformation in $\mathsf{CSim}_4$.

Each coefficient of $y_0$ has parity bias $\delta_0(r)$. Each folded bit sums two independent coefficients, so its parity bias is $\delta_0(r)^2$. The $b$ folded bits are independent. Every fixed folded vector therefore has probability at most $((1+\delta_0(r)^2)/2)^b$. Both commitments contain this vector, proving~\cref{eq:fold-entropy}.
\end{proof}

\subsection{Four-Move Unforgeability}

\begin{proof}[Proof of~\cref{thm:unit-security}]
Correctness is~\cref{lem:unit-correctness,lem:four-compressed}. For an adversary with no signing queries, replace the public key by uniform $(\mat A_0,\vec b)$ under the bounded-key NTWE assumption. Every basic forgery with nonzero response solves the folded self-target problem at bound $B_s$. Every compressed forgery determines a response $\vec z'$ with $\mat B\vec z'=\vec w_h'\pmod q$ and unchanged first ring coordinate. Thus a nonzero reconstructed response solves the same problem at bound $B_v$.

A zero response requires $H_4(\vec0,c,\mu)=c$. Distinct challenges give distinct oracle inputs, so the contribution of zero responses is at most $(q_H+1)/|\mathcal C_4|$, including an unqueried final output. This completes the no-message argument. Apply~\cref{lem:fs-transfer} with $\gamma=\gamma_4$ and $\beta=1-\eta M^{-\kappa}$. Its simulation and entropy hypotheses follow from~\cref{lem:four-simulation}, and its additional terms are negligible under the stated conditions. Total negligible implementation error changes the success probability negligibly. For either two-move rule, $\pm1=1\pmod I$ preserves correctness, and its incoming identity gives the same simulation argument.
\end{proof}
\FloatBarrier

% !TEX root = main.tex
\section{Gaussian Sum Identities and Validation}
\label{app:analytic}

\subsection{The Largest Gaussian Parity Imbalance}

\begin{proof}[Proof of~\cref{lem:worst}]
The function $\Te$ is even and $2$-periodic, and $\To(t)=\Te(t+1)$. Write $Q=e^{-\pi^2\alpha^2/2}$, so $0<Q<1$. Poisson summation for the Gaussian and the theta product identity~\cite[Equation~20.5.3]{DLMF} give the following series and product for the same function,
\begin{align}
 \Te(t)
 &=\frac{\alpha\sqrt{2\pi}}2\left(1+2\sum_{m=1}^{\infty}Q^{m^2}\cos(\pi mt)\right)\label{eq:fourier}\\
 &=\frac{\alpha\sqrt{2\pi}}2\prod_{m=1}^{\infty}(1-Q^{2m})
 \bigl(1+2Q^{2m-1}\cos(\pi t)+Q^{4m-2}\bigr).\label{eq:product}
\end{align}
The Fourier coefficient of $e^{\pi i mt}$ in the periodised Gaussian is
\[
 \frac12\int_{\R}e^{-u^2/(2\alpha^2)}e^{-\pi imu}\,du=\frac{\alpha\sqrt{2\pi}}2Q^{m^2},
\]
which gives~\cref{eq:fourier}. The standard identity used in~\cref{eq:product} is
\[
 1+2\sum_{m\ge1}Q^{m^2}\cos(2mz)
 =\prod_{m\ge1}(1-Q^{2m})(1+2Q^{2m-1}\cos(2z)+Q^{4m-2}),
\]
with $z=\pi t/2$.

Every factor in the second parentheses in~\cref{eq:product} is positive. For $0\le s<t\le1$, each such factor at $s$ is strictly larger than at $t$, because $\cos(\pi s)>\cos(\pi t)$. The products converge to positive limits since the sums of $Q^{2m}$ and $Q^{2m-1}$ converge. The ratio of the products at $s$ and $t$ is therefore strictly larger than one, already because its first factor is larger than one and all subsequent factors are larger than one. Hence $\Te$ strictly decreases on $[0,1]$.

On $[0,1]$, evenness and periodicity give $\To(t)=\Te(1-t)$. Therefore $\Te(t)/\To(t)$ strictly decreases from $\Te(0)/\Te(1)$ to its reciprocal. Its value at $t=1/2$ is one. Periodicity reduces all other $t$ to this interval, possibly after exchanging the two classes. This proves the maximum and also $\Te(0)>\Te(1)$.

Finally, evaluating~\cref{eq:product} at $t=0$ and $t=1$ gives
\begin{equation}
\label{eq:rate-product}
 M^*(\alpha)=\prod_{m=1}^{\infty}\left(\frac{1+Q^{2m-1}}{1-Q^{2m-1}}\right)^2.
\end{equation}
Each factor strictly increases with $Q$, and $Q$ strictly decreases with $\alpha$. Comparing the convergent positive products, as above, proves the final claim.
\end{proof}

\subsection{Strict Improvement over the Original Functions}

\begin{proof}[Proof of~\cref{prop:improvement}]
Set $\Delta=\Te(0)-\To(0)>0$ and $p_1=e^{-1/(2\alpha^2)}$. Since $\To(0)>p_1$, we have $M^*(\alpha)=1+\Delta/\To(0)<1+\Delta/p_1$. Subtracting~\cref{eq:fourier} at $t=1$ from its value at $t=0$ gives, for $Q=e^{-\pi^2\alpha^2/2}$,
\[
 \Delta=2\alpha\sqrt{2\pi}\sum_{j\ge0}Q^{(2j+1)^2}
 \le\frac{2\alpha\sqrt{2\pi}\,Q}{1-Q^4}.
\]
The inequality follows from $(2j+1)^2\ge1+4j$ and summing a geometric series. Substituting this bound into $1+\Delta/p_1$ gives~\cref{eq:original}.

For the second claim, the original functions at input $3\vec v$ have probabilities $S/M$ and $(1-S)/M$, where
\[
 S=\sum_{j\ge0}(-1)^j\frac{\rho_r((3+j)\vec v)}{\rho_r(3\vec v)}.
\]
Thus this odd input accepts with probability $1/M$. Set $p_3=e^{-9/(2\alpha^2)}$. Summing~\cref{eq:balance} over the even points of $\Z\vec v$, and bounding the acceptance probabilities of all other odd inputs by one, gives
\[
 \frac{\Te(0)}M\le\To(0)-p_3+\frac{p_3}M.
\]
Since $0<p_3<\To(0)$ and $\Te(0)>\To(0)$, it follows that
\[
 M\ge\frac{\Te(0)-p_3}{\To(0)-p_3}>\frac{\Te(0)}{\To(0)}=M^*(\alpha).
\]
\end{proof}

\subsection{Numerical Comparisons and Checks}

The comparison in~\cref{tab:rates} evaluates the exact optimum of~\cref{thm:optimal} against the explicit factor of~\cref{def:original}. The latter is the sufficient factor used in~\cite{GartnerFull}, rather than a new optimisation of its original acceptance functions. A reduction in per-step rejection measures how often this step aborts, and does not directly give a signing-time speedup.

\begin{table}[H]
\centering
\begin{tabular}{cccc}
\toprule
$\alpha$ & $1-1/\Mold(\alpha)$ & $1-1/M^*(\alpha)$ & Rejection reduction\\
\midrule
$0.8$ & $2^{-1.88}$ & $2^{-2.67}$ & $42.2\%$\\
$1.0$ & $2^{-4.16}$ & $2^{-5.14}$ & $49.5\%$\\
$1.2$ & $2^{-7.17}$ & $2^{-8.25}$ & $52.8\%$\\
$1.5$ & $2^{-12.79}$ & $2^{-14.02}$ & $57.4\%$\\
\bottomrule
\end{tabular}
\caption{Per-step rejection probabilities, rounded.}
\label{tab:rates}
\end{table}

The reduction is $1-(1-1/M^*)/(1-1/\Mold)$. We also evaluated the functions on $170$ Gaussian lines, using $\alpha\in\{0.15,0.25,0.4,0.6,0.8,1,1.2,1.5,2,4\}$ and offsets $t=i/16$ for $i=0,\ldots,16$. Across $7752$ point and repetition-factor combinations, both probabilities were nonnegative and their sum was at most one within error $2^{-270}$. The relative residual in~\cref{eq:balance} and the absolute residual in~\cref{eq:sum-budget} were below $2^{-270}$. The checks used the optimal factor for each line, the common factor $M^*(\alpha)$, and $2M^*(\alpha)$.

The mathematical statements concern the infinite sums with exact real arithmetic. A finite evaluation must control the error in their ratios as well as the discarded Gaussian tails.
The four-move acceptance identities were checked on $96$ orthogonal point and factor combinations and $20$ ring-unit instances, at width ratios $0.4$, $0.8$, $1.2$, and $2$. Relative residuals were below $2^{-300}$. Compression was checked exhaustively on $62534$ scalar residues, including all residues at the two concrete moduli, and on $64$ ring transcripts. An additional $192$ checks rejected malformed hints, challenges outside the prescribed support, and responses exceeding the verification bound.
\FloatBarrier

% !TEX root = main.tex
\section{Detailed Two-Move Parameters and Estimates}
\label{sec:comparison}

The new functions can reduce rejection at the existing Gaussian width or reduce the width at the existing repetition factor. We evaluate both choices for the three parameter sets of~\cite[Table~1]{GartnerFull}. The parameters held fixed are given in~\cref{tab:parameters}. Set names identify parameter sets, not newly established security levels.

Both the original and improved two-move columns use binary challenges with the same reduced Hamming weights. Sets I and II use weights exactly $58$ and $80$. Set III uses maximum weight $128$ and the half-space choice defined after~\cref{def:parameters}, with $2^{255}$ challenges. None of these three sets uses the full binary challenge space. The exponent $\kappa$ in the rejection calculation is the fixed weight or the maximum permitted weight, never an average weight.

\begin{definition}[Coefficient-Level Lattice Parameters]
\label{def:lattice-dimensions}
For ring degree $d$ and module parameters $\ell,m$, the verification matrix has $m$ rows and $\ell+m+1$ columns over the ring. Its coefficient matrix has
\[
 n_{\mathsf{SIS}}\coloneqq dm,\qquad
 m_{\mathsf{SIS}}\coloneqq d(\ell+m+1)=N.
\]
The unstructured SIS proxy uses this shape, modulus $q$, and Euclidean bound $B_v$, or $2B_v$ for the response-difference diagnostic. The NTWE proxy of~\cite[Section~5.1 and Appendix~E]{GartnerFull} is unstructured LWE with
\[
 n_{\mathsf{LWE}}\coloneqq d(\ell+1)-1,\qquad
 X_s=X_e=D_{\Z,\sigma},\qquad m_{\mathsf{LWE}}=\infty.
\]
The actual module instance supplies $dm$ scalar equations. The infinite-sample proxy is the convention used for the cited estimates.
\end{definition}

The integer SIS lattice has rank $N$, not $N-dm$, because it contains $q\Z^N$. The verification matrix modulo $q$ contains an invertible scaled identity block, so its kernel over $\ring_q$ is a free module of rank $\ell+1$. These are different uses of rank.

\cref{tab:parameters} reproduces the rejecting parameters and security estimates of~\cite[Table~1]{GartnerFull}. In each security cell, the pair is the BKZ block size followed by the classical Core-SVP exponent $\lfloor0.292\beta\rfloor$. All nine pairs reproduce exactly using the unstructured instances listed in Appendix~E there. These estimates are heuristic work factors, not bounds on the full signature-forgery advantage.

\begin{table}[!htbp]
\centering
\begin{tabular}{lrrr}
\toprule
Parameter & I & II & III\\
\midrule
Ring degree $d$ & 256 & 256 & 256\\
Modulus $q$ & 12289 & 50177 & 50177\\
Module parameters $(\ell,m)$ & $(1,2)$ & $(3,2)$ & $(4,3)$\\
Module matrix shape & $2\times4$ & $2\times6$ & $3\times8$\\
Coefficient matrix shape & $512\times1024$ & $512\times1536$ & $768\times2048$\\
NTWE proxy dimension & 511 & 1023 & 1279\\
Secret width $\sigma$ & 2.60 & 1.00 & 1.50\\
Response width $r$ & 128 & 55 & 95\\
Maximum challenge weight $\kappa$ & 58 & 80 & 128\\
Compression parameter $\tau$ & 256 & 128 & 128\\
Secret bound $B_k$ & 110.07 & 48.00 & 79.60\\
Signing bound $B_s$ & 4178 & 2199 & 4386\\
Listed verification bound $B_v$ & 5649 & 2946 & 5300\\
Key bytes & 928 & 1056 & 1568\\
Signature bytes & 775 & 1184 & 1694\\
SIS at $B_v$, block size / bits & $420/122$ & $629/183$ & $899/262$\\
SIS at $2B_v$, block size / bits & $332/96$ & $500/146$ & $730/213$\\
NTWE, block size / bits & $414/120$ & $618/180$ & $881/257$\\
\bottomrule
\end{tabular}
\caption{The rejecting parameter sets and estimates of~\cite[Table~1]{GartnerFull}.}
\label{tab:parameters}
\end{table}

The listed value $B_v=5300$ in set III is retained for exact reproduction of the security estimates. Enforcing~\cref{eq:verification-bound} with the displayed integer $B_s=4386$ requires $5301$, as quantified in~\cref{tab:bounds}. This adjustment does not change the signature-size estimate.

\subsection{Fewer Rejection Attempts at Fixed Width}

\begin{proposition}[Rejection Comparison]
\label{prop:concrete-rejection}
For independent uniform challenges in the identification protocol, the rejection probability of one step is $1-M^{-1}$, and the rejection probability of the iterative sampler, including its final retention test, is $1-M^{-\kappa}$. The expected number of independent sampler attempts until acceptance is $M^\kappa$. At fixed $r$ and $B_k$, replacing the original functions by the improved functions therefore reduces this expectation by the fraction
\[
 1-\left(\frac{M^*(r/B_k)}{\Mold(r/B_k)}\right)^\kappa.
\]
\end{proposition}
\begin{proof}
The step probability follows from~\cref{lem:output} and the sampler probability from~\cref{cor:iteration}. Independent attempts have a geometric waiting time with mean equal to the inverse success probability. Substitution of the two repetition factors gives the reduction.
\end{proof}

\cref{tab:fixed-width} evaluates~\cref{prop:concrete-rejection} with the printed $B_k$ values of~\cref{tab:parameters}. These probabilities describe the sampler stage before the response-norm test or a byte-encoding overflow test. They also give the signing-attempt model with fresh independent random-oracle challenges. They are not timings of an implementation. At fixed width, the key sizes, response bounds, and signature-size model remain unchanged.

\begin{table}[t]
\centering
\begin{tabular}{lrrrrrr}
\toprule
& & \multicolumn{2}{c}{Sampler rejection} & \multicolumn{2}{c}{Expected attempts} & Reduction\\
Set & $r$ & Original & Improved & Original & Improved & \\
\midrule
I & 128 & $2^{-1.122}$ & $2^{-1.976}$ & 1.850 & 1.341 & $27.5\%$\\
II & 55 & $2^{-0.640}$ & $2^{-1.365}$ & 2.790 & 1.634 & $41.4\%$\\
III & 95 & $2^{-0.696}$ & $2^{-1.455}$ & 2.612 & 1.574 & $39.7\%$\\
\bottomrule
\end{tabular}
\caption{Before and after at the original Gaussian widths.}
\label{tab:fixed-width}
\end{table}

\subsection{Smaller Widths at Fixed Rejection}

\begin{definition}[Matched Width and Size Model]
\label{def:size-model}
For an original width $r_{\mathsf{orig}}$, define $r_{\mathsf{new}}$ by
\begin{equation}
\label{eq:matched-rate}
 M^*(r_{\mathsf{new}}/B_k)=\Mold(r_{\mathsf{orig}}/B_k).
\end{equation}
Following~\cite[Equation~(8) and Appendix~A.2]{GartnerFull}, define the compressed-signature entropy estimate in bits and the corresponding padded byte estimate by
\begin{align}
 H_{\mathsf{est}}(r)&\coloneqq\frac{d(\ell+1)}2\log_2(2\pi e r^2)+\frac{dm}2\log_2\bigl(8\pi e(r/\tau)^2\bigr),\label{eq:entropy-size}\\
 L_{\mathsf{est}}(r)&\coloneqq\left\lceil H_{\mathsf{est}}(r)/8\right\rceil+64.\label{eq:byte-size}
\end{align}
The additive $64$ bytes account for a $256$-bit challenge and $32$ bytes of padding in this model. The public-key estimate is $32+\lceil dm\lceil\log_2q\rceil/8\rceil$ bytes, using a $32$-byte seed for $\mat A_0$. This seed-based estimate presumes pseudorandom matrix expansion, whereas the mathematical protocol samples $\mat A_0$ uniformly.
\end{definition}

\begin{proposition}[Matched-Rate Size Difference]
\label{prop:size-difference}
The matched width exists uniquely and satisfies $0<r_{\mathsf{new}}<r_{\mathsf{orig}}$. It gives the same iterative-sampler acceptance probability as the original functions at $r_{\mathsf{orig}}$. For fixed $d,\ell,m,\tau$, the entropy difference is
\[
 H_{\mathsf{est}}(r_{\mathsf{orig}})-H_{\mathsf{est}}(r_{\mathsf{new}})
 =N\log_2(r_{\mathsf{orig}}/r_{\mathsf{new}}).
\]
\end{proposition}
\begin{proof}
By~\cref{lem:worst}, $M^*(\alpha)$ is continuous and strictly decreasing. Its defining sums give $M^*(\alpha)\to\infty$ as $\alpha\to0$, and its product formula gives $M^*(\alpha)\to1$ as $\alpha\to\infty$. Hence~\cref{eq:matched-rate} has a unique solution. By~\cref{prop:improvement}, $M^*(r_{\mathsf{orig}}/B_k)<\Mold(r_{\mathsf{orig}}/B_k)$, so the solution has smaller width. Equality of the factors gives equality of $M^{-\kappa}$. Finally, the coefficients of $\log_2r$ in~\cref{eq:entropy-size} sum to $d(\ell+1)+dm=N$.
\end{proof}

\cref{tab:fixed-rate} gives the matched widths and size estimates. Each byte count is rounded separately according to~\cref{eq:byte-size}. The widths shown in the table are rounded for display. The estimates reproduce the original $775$, $1184$, and $1694$ bytes and become $764$, $1167$, and $1674$ bytes. Actual encoded lengths and encoding-induced rejection depend on the chosen entropy coder.

\begin{table}[t]
\centering
\begin{tabular}{lrrrrrr}
\toprule
& \multicolumn{2}{c}{Gaussian width} & \multicolumn{2}{c}{Estimated signature bytes} & \multicolumn{2}{c}{Saved}\\
Set & Original & Improved & Original & Improved & Bytes & Fraction\\
\midrule
I & 128 & 120.6834 & 775 & 764 & 11 & $1.42\%$\\
II & 55 & 51.7805 & 1184 & 1167 & 17 & $1.44\%$\\
III & 95 & 89.7892 & 1694 & 1674 & 20 & $1.18\%$\\
\bottomrule
\end{tabular}
\caption{Before and after at the original iterative-sampler rejection rates.}
\label{tab:fixed-rate}
\end{table}

The norm bounds in~\cref{tab:bounds} use $B_s=\lceil1.02r\sqrt N\rceil$ and $B_v=\lceil B_s+\sqrt{dm}(\tau/4+1)\rceil$. For set III at the original width, this gives $B_v=5301$, whereas~\cite[Table~1]{GartnerFull} lists $5300$. The value $5301$ satisfies~\cref{eq:verification-bound}, since $4386+33\sqrt{768}>5300$. This rounding adjustment does not affect the rejection or size calculations.

\begin{table}[t]
\centering
\begin{tabular}{lrrrr}
\toprule
& \multicolumn{2}{c}{Original width} & \multicolumn{2}{c}{Matched width}\\
Set & $B_s$ & $B_v$ & $B_s$ & $B_v$\\
\midrule
I & 4178 & 5649 & 3940 & 5411\\
II & 2199 & 2946 & 2070 & 2817\\
III & 4386 & 5301 & 4145 & 5060\\
\bottomrule
\end{tabular}
\caption{Signing and verification bounds for the two width choices.}
\label{tab:bounds}
\end{table}

The matched-width comparison preserves the key distribution, challenge space, and modulus, and reduces the response bound in the conditional self-target assumption of~\cref{thm:signature-security}. Its fixed-rate condition concerns iterative rejection. The separate norm and encoding tests must still be included when measuring full signing performance.

\FloatBarrier
% !TEX root = main.tex
\section{Detailed Four-Move Parameters and Estimates}
\label{app:four-comparison}

We compare three distinct lattice parameter sets. Within each set, the key distribution, challenge space, and compressed verification relation are fixed. The three columns use the two-move rule of~\cite{GartnerFull}, optimal two moves, and optimal four moves. The two-move rules are valid in the folded protocol because $\pm1=1\pmod I$. The columns therefore compare acceptance rules on the same parameters. The rejecting parameter sets of~\cite{GartnerFull} are reproduced separately in~\cref{tab:parameters}.

\begin{definition}[Classical Query Budget]
\label{def:query-budget}
Let $Q$ be the maximum number of classical queries made to $H_4$ in the no-message challenge-search calculation, and put $h\coloneqq\log_2|\mathcal C_4|$. The illustrative budget is $Q=2^{64}$. A target exponent $t$ allocates $2^{-t}$ to the query contribution $Q2^{-h}$, requiring $h\ge t+\log_2Q$. The query count $Q$ is distinct from the modulus $q$.
\end{definition}

\begin{definition}[Concrete Comparison Parameters]
\label{def:four-parameters}
Use sets A, B, and C of~\cref{tab:four-parameters}, all with
\[
 q=50177,\qquad \Delta=64,\qquad
 B_k=\left\lceil\sigma\sqrt{d(\ell+m+4)}\right\rceil.
\]
The binary challenges are supported on the first $b=d/2$ coefficients. Their distributions are
\begin{equation}
\label{eq:four-weight}
 \begin{array}{c|c|c}
 \text{Set}&\mathcal C_4&|\mathcal C_4|\\ \hline
 A&\{c\in\{0,1\}^{256}:\norm{c}_1=58\}&\binom{256}{58}\\
 B&\{0,1\}^{256}&2^{256}\\
 C&\{c\in\{0,1\}^{512}:\norm{c}_1=82\}&\binom{512}{82}
 \end{array}
\end{equation}
For each set, choose widths $r_{\mathsf{orig}},r_2,r_4$ with two expected sampler attempts,
\begin{equation}
\label{eq:four-matched}
 \Mold(r_{\mathsf{orig}}/B_k)^\kappa
 =M^*(r_2/B_k)^\kappa
 =M_4^*(r_4/B_k)^\kappa=T,\qquad T=2.
\end{equation}
\end{definition}


Sets A and C use reduced, fixed Hamming weights. Set B uses weights from $0$ to $256$, with mean $128$. Its repetition exponent is the maximum $256$, not the mean. The final retention test makes acceptance independent of the actual weight. Each row uses the same challenge distribution for both two-move rules and four moves.

The target exponent $t$ in~\cref{tab:four-parameters} concerns only challenge search. It does not label a complete signature-security level. At degree $512$, the smaller target in set A permits weight $58$. Set B needs all $256$ challenge bits for its allocation. Set C uses degree $1024$ to obtain enough challenge positions for its larger target while retaining a fixed weight of $82$.

\subsection{Smaller Signatures at Fixed Rejection}

\begin{definition}[Four-Move Size Estimate]
\label{def:four-size}
For each parameter set, define the challenge encoding length
\[
 L_c\coloneqq\max\left\{32,\left\lceil\frac{\log_2|\mathcal C_4|}{8}\right\rceil\right\}
\]
and the estimates
\begin{align}
 H_{4,\mathsf{est}}(r)&\coloneqq
 \frac{d(\ell+1)}2\log_2(2\pi e r^2)
 +\frac{dm}2\log_2\bigl(2\pi e(r/\Delta)^2\bigr),\label{eq:four-entropy-size}\\
 L_{4,\mathsf{est}}(r)&\coloneqq
 \left\lceil H_{4,\mathsf{est}}(r)/8\right\rceil+L_c+32.\label{eq:four-byte-size}
\end{align}
The final $32$ bytes are padding. The response coordinates use the Gaussian entropy model of~\cref{def:size-model}, and the hint indices use Gaussian width $r/\Delta$.
\end{definition}

A fixed-weight challenge can be represented by its index in a fixed enumeration of the corresponding binary vectors. Thus $L_c$ bytes suffice for an injective challenge encoding. Sets A and B retain the $32$ challenge bytes used in~\cite{GartnerFull}. Set C requires $41$ bytes. The same overhead is included in every sampler column of a row.

Setting $\Delta=\tau/2$ makes~\cref{eq:four-entropy-size} equal to~\cref{eq:entropy-size}. The coefficient of the omitted response is one instead of two, giving hint width $r/\Delta$ and worst-case reconstruction error $\Delta/2+1=\tau/4+1$. The byte counts are entropy estimates. A concrete encoder must also reject malformed encodings and account for overflow.

\begin{proposition}[Fixed-Rate Four-Move Improvement]
\label{prop:four-size}
At the matched factors of~\cref{eq:four-matched}, $r_4<r_2<r_{\mathsf{orig}}$, and the entropy saving relative to either two-move width $r_i$ is
\[
 H_{4,\mathsf{est}}(r_i)-H_{4,\mathsf{est}}(r_4)
 =N\log_2(r_i/r_4).
\]
All three samplers have acceptance probability $1/T$ before the norm test.
\end{proposition}
\begin{proof}
For every $\alpha>0$, $1<M_4^*(\alpha)<M^*(\alpha)<\Mold(\alpha)$. The first comparison follows from~\cref{eq:four-rate}, and the last is~\cref{prop:improvement}. The optimal factors are continuous and strictly decreasing, with limits infinity and one. Hence matching a common factor greater than one gives the stated width ordering. The coefficients of $\log_2r$ in~\cref{eq:four-entropy-size} sum to $N$, proving the saving. The acceptance claim follows from the iterative incoming identities and the retention test.
\end{proof}

\cref{tab:four-sizes} gives the size comparison at two expected sampler attempts in every row. Four moves save $162$, $195$, and $259$ bytes relative to the two-move rule of~\cite{GartnerFull} on the same parameters. Optimising the two-move functions accounts for $28$, $28$, and $42$ bytes, while the additional orthogonal moves account for another $134$, $167$, and $217$ bytes. The key size does not change between columns.


\subsection{Widths, Norm Bounds, and Lattice Estimates}

\cref{tab:four-bounds} reports every response width and norm bound used in~\cref{tab:four-sizes}. The signing cutoff is $B_s=\lceil1.02r\sqrt N\rceil$, and the verification cutoff is $B_v=\lceil B_s+\sqrt{dm}(\Delta/2+1)\rceil$. The comparison uses the same folded relation within each set, so~\cref{thm:unit-security} applies to all three rules under the folded self-target assumption at that row's largest verification bound.


\cref{tab:four-lattice} evaluates the unstructured proxies of~\cref{def:lattice-dimensions} at the matrix dimensions of~\cref{tab:four-parameters}. Each pair again gives the BKZ block size and classical Core-SVP exponent. The NTWE estimates are $969/283$ for set A and $1383/404$ for both B and C, unchanged between sampler columns. Their equality for B and C is a property of the shared unstructured LWE proxy parameters, not an assertion that the different ring structures have identical security.

\begin{table}[!htbp]
\centering
\begin{tabular}{llrr}
\toprule
Set & Acceptance rule & SIS at $B_v$ & SIS at $2B_v$\\
\midrule
A & Rule of~\cite{GartnerFull} & $1379/402$ & $1121/327$\\
  & Optimal two & $1399/408$ & $1136/331$\\
  & Four & $1499/437$ & $1208/352$\\
\midrule
B & Rule of~\cite{GartnerFull} & $1265/369$ & $1035/302$\\
  & Optimal two & $1281/374$ & $1047/305$\\
  & Four & $1377/402$ & $1121/327$\\
\midrule
C & Rule of~\cite{GartnerFull} & $2947/860$ & $2364/690$\\
  & Optimal two & $2994/874$ & $2396/699$\\
  & Four & $3253/949$ & $2569/750$\\
\bottomrule
\end{tabular}
\caption{Unstructured SIS proxies, given as BKZ block size / classical Core-SVP bits.}
\label{tab:four-lattice}
\end{table}

These lattice estimates do not establish the hardness of the folded self-target relation. In particular, they do not account for challenge search. The $2B_v$ column records the norm bound on the difference of two accepted responses and is not a strong-unforgeability theorem for the four-move scheme. The security statement remains the conditional ordinary-unforgeability theorem of~\cref{thm:unit-security}.

\subsection{Lower Rejection at Fixed Signature Size}

Alternatively, retain $r_{\mathsf{orig}}$ from~\cref{tab:four-bounds} within each set. All three rules then have the same signature-size estimate and norm bounds. \cref{tab:four-fixed-width} reports their sampler rejection $1-M^{-\kappa}$ and expected attempts $M^\kappa$. These are repetition counts, not running-time measurements, since the four-move step evaluates an additional direction.

\begin{table}[!htbp]
\centering
\begin{tabular}{lrrrrrr}
\toprule
& & \multicolumn{2}{c}{Sampler rejection} & \multicolumn{3}{c}{Expected attempts}\\
Set & Bytes & Optimal two & Four & Rule of~\cite{GartnerFull} & Two & Four\\
\midrule
A & 1892 & $2^{-1.826}$ & $2^{-10.046}$ & 2 & 1.392827 & 1.000947\\
B & 2557 & $2^{-1.873}$ & $2^{-12.297}$ & 2 & 1.375682 & 1.000199\\
C & 2991 & $2^{-1.837}$ & $2^{-10.571}$ & 2 & 1.388798 & 1.000658\\
\bottomrule
\end{tabular}
\caption{Fixed-size rejection, with original-rule rejection $2^{-1}$ in every row.}
\label{tab:four-fixed-width}
\end{table}

\subsection{Norms Used in the Size Calculation}

The public challenge has coefficients in $\{0,1\}$, not a centred alphabet. A nonzero bit requests a step, while the acceptance rule chooses its sign and direction. Thus the response is $\vec z=\vec y+c'\vec s$, not necessarily $\vec y+c\vec s$. The Gaussian response law follows exactly from the incoming identities, as in~\cref{lem:four-simulation}. It does not follow from a central-limit approximation or an assumption of independent random signs.

Each step shifts by $X^j\vec s$ or $\omega X^j\vec s$, both of norm at most $B_k$. The raw sum satisfies
\[
 \norm{c'\vec s}\le\norm{c'}_1\norm{\vec s}
 =k\norm{\vec s}\le\kappa B_k,
\]
because the lifted challenge has one signed coefficient for every nonzero challenge bit. This worst-case bound is not used to set the response width. A heuristic $\sqrt{k}$ expansion is not used either. The size model uses the accepted Gaussian response law, while every accepted signature satisfies the hard verification cutoff and the compression error is bounded in the worst case.

\begin{lemma}[Gaussian Norm Tail, {\cite[Lemma~6]{GartnerFull}}]
\label{lem:multi-norm-tail}
For $\vec z\leftarrow D_{\Z^N,r}$ and $C\ge1$,
\[
 \Pr[\norm{\vec z}>Cr\sqrt N]\le
 \left(C\exp\!\left(\frac{1-C^2}{2}\right)\right)^N.
\]
\end{lemma}

The norm test multiplies the sampler acceptance by $\eta$ from~\cref{lem:four-simulation}. Thus $T$ describes sampler attempts, while full signing also includes norm and encoding rejection. The multiplier $1.02$ is a cutoff near the typical Gaussian norm, not a claim that norm rejection is at most $2^{-192}$. Folding changes only commitment parity, not the dimension or norm of the transmitted response.

\subsection{The Hash-Query Budget}



\begin{proposition}[Challenge Search]
\label{prop:challenge-search}
For $Q\le|\mathcal C_4|$ distinct challenges and one unsigned message, a classical adversary can produce a zero-response forgery with probability
\[
 1-(1-2^{-h})^Q.
\]
The zero-response contribution of any classical adversary with at most $Q$ oracle queries is at most $(Q+1)2^{-h}$.
\end{proposition}
\begin{proof}
For each tested challenge $c$, query $H_4(\vec0,c,\mu)$. If the value equals $c$, the basic signature $(\vec0,c)$ and the compressed signature $(\vec0,\vec0,c)$ both verify. Distinct challenges give independent oracle inputs, each with success probability $2^{-h}$. This proves the first formula. For the upper bound, each query can match at most one predetermined challenge. Apply the union bound and include a possible unqueried final output.
\end{proof}

The challenge entropies and query contributions for all three sets are given in~\cref{tab:four-parameters}. In set B, $Q2^{-h}=2^{64-256}=2^{-192}$, with an additional $2^{-256}$ for a possible unqueried final output. Sets A and C have enough entropy slack to include this final output within their respective $2^{-128}$ and $2^{-256}$ allocations. A reduced weight at $b=256$ cannot attain the full $256$ challenge bits required by set B's query allocation. Increasing the degree changes this restriction, as the fixed-weight choice in set C demonstrates.

The full chosen-message bound also includes the advantages against the folded self-target and bounded-key NTWE assumptions, and the simulation terms of~\cref{lem:fs-transfer}. The query calculation alone does not establish a total $2^{-192}$ forgery bound. A quantum query budget has a different search cost~\cite{GroverSearch,BBHTSearch}, and a success bound is distinct from the work-factor comparison used in NIST security categories~\cite{NISTCategories}.

\FloatBarrier

\end{document}
